%%%%%%%%%%%%%%%%%%%%%%%%%%%%%%%%%%%%%%%%%%%%%%%%%%%%%%%%%%%%%%%%%%%%%%%%%
%  Deep Learning-Based Cocoa Pod Disease Detection                     %
%  Extended IEEE Conference Format Paper (15-20 pages with appendix)   %
%  COE 557 - Advanced Artificial Intelligence                          %
%%%%%%%%%%%%%%%%%%%%%%%%%%%%%%%%%%%%%%%%%%%%%%%%%%%%%%%%%%%%%%%%%%%%%%%%%

\documentclass[conference,10pt]{IEEEtran}
\IEEEoverridecommandlockouts

\usepackage{cite}
\usepackage{amsmath,amssymb,amsfonts}
\usepackage{algorithmic}
\usepackage{graphicx}
\usepackage{textcomp}
\usepackage{xcolor}
\usepackage{booktabs}
\usepackage{multirow}
\usepackage{array}
\usepackage{subcaption}
\usepackage{url}
\usepackage[hidelinks]{hyperref}
\usepackage{listings}
\usepackage{diagbox}
\usepackage{placeins}
\usepackage{stfloats}

% so that figures in paper/Images can be found by Overleaf and locally
\graphicspath{{../Images/}{Images/}{./}}

\def\BibTeX{{\rm B\kern-.05em{\sc i\kern-.025em b}\kern-.08em
    T\kern-.1667em\lower.7ex\hbox{E}\kern-.125emX}}

\begin{document}

\title{Deep Learning-Based Cocoa Pod Disease Detection:\\
A Multi-Class Classification System for Ghanaian Farmers}

\author{
\IEEEauthorblockN{Enoch K. Koranteng}
\IEEEauthorblockA{
Department of Computer Engineering\\
Kwame Nkrumah University of Science and Technology\\
Kumasi, Ghana\\
ekkoranteng10@st.knust.edu.gh
}
}

\maketitle

\begin{abstract}
Cocoa (\textit{Theobroma cacao} L.) is a cornerstone of Ghana's economy,
generating approximately US\$2 billion in annual foreign exchange and supporting
over 800{,}000 farm families. Yet fungal and pest diseases, including
\textit{Phytophthora palmivora} (black pod), \textit{Moniliophthora roreri}
(frosty pod rot), \textit{Moniliophthora perniciosa} (witches' broom), and
\textit{Carmenta foraseminis} (cocoa pod borer), destroy 30--40\% of global
cocoa production each year. Early identification by trained agronomists is
chronically scarce in rural Ghana. This paper presents a deep-learning
multi-class classification system that identifies five pod conditions
(including a healthy class) directly from smartphone images. Three CNN
architectures, ResNet-50, EfficientNet-B0, and MobileNetV3-Small, are
adapted via ImageNet transfer learning and fine-tuned on a curated dataset of
3{,}285 annotated pod crops. To address severe class imbalance (the rarest
class is 2.3\% of samples), we apply inverse-frequency cross-entropy weighting.
Evaluation covers accuracy, macro F1, per-class precision/recall, ROC AUC,
inference latency distribution, parameter count, and FP32/INT8 model footprint
to assess both predictive quality and edge feasibility. ResNet-50 achieves the
highest test accuracy (87.25\%), EfficientNet-B0 the best macro F1 (0.755) at
4.01M parameters, and MobileNetV3-Small the lowest latency (3.9 ms) at 1.52M
parameters and an estimated 1.45 MB INT8 footprint, satisfying the constraints
of offline smartphone deployment. The contribution is a complete, reproducible
pipeline plus an analysis of where each architecture sits on the
accuracy--efficiency frontier for Ghanaian smallholder farms.
\end{abstract}

\begin{IEEEkeywords}
Cocoa pod disease, deep learning, convolutional neural networks, transfer
learning, ResNet-50, EfficientNet, MobileNetV3, precision agriculture, Ghana,
edge deployment, ROC analysis, model compression.
\end{IEEEkeywords}

%==============================================================
\section{Introduction}\label{sec:intro}

\IEEEPARstart{G}{hana} is the world's second-largest producer of cocoa, contributing
approximately 20\% of global output \cite{usda2025}. The crop underpins the
livelihoods of nearly 800{,}000 farm families and is the country's third-largest
source of foreign exchange revenue, surpassed only by gold and crude oil
\cite{gcb2022}. Despite this economic significance, cocoa yields in Ghana have
declined sharply: the 2024/2025 crop season saw output nearly halve compared
to the 2021/2022 peak of 1.04 million metric tons, with disease and climate
stress cited as primary drivers \cite{usda2025}.

The principal pathogens threatening Ghanaian cocoa are well documented.
\textit{Phytophthora palmivora}, the causal agent of black pod disease,
accounts for 10--30\% of annual pod losses globally and can destroy up to 90\%
of a crop if left untreated \cite{alvarado2025}.
\textit{Moniliophthora roreri} (frosty pod rot) and \textit{Moniliophthora
perniciosa} (witches' broom) are predominantly a threat in Latin America but
have been detected in West African farms and represent a biosecurity risk
\cite{buenano2024}. \textit{Carmenta foraseminis}, the cocoa pod borer,
inflicts yield losses of up to 75\% in severely affected areas
\cite{pagaduan2026}. Cocoa swollen-shoot virus disease (CSSVD) compounds these
threats; over 90{,}000 hectares of Ghanaian cocoa are currently infected, and
more than 200 million trees have been felled since the disease was first
documented in 1936 \cite{gcb2022}.

\subsection{Why Smartphone-Based Diagnosis}

Conventional disease diagnosis requires field visits from trained agricultural
extension officers, a resource that is chronically scarce in rural Ghana.
Farmers typically wait days or weeks for a diagnosis, by which time disease
spread has become difficult to contain. Smartphone penetration in sub-Saharan
Africa has grown rapidly, reaching over 50\% in Ghana by 2023, creating an
opportunity for AI-powered on-device diagnostics that are accessible, fast,
and low-cost \cite{cocoacompanion2022}. The deployment context imposes three
constraints simultaneously: (i)~the model must run \emph{offline}, because
many cocoa-growing regions have intermittent or no mobile data; (ii)~latency
must be perceptually instant ($<\!100$\,ms) on a 2--4 GB Android device; and
(iii)~the binary footprint must fit alongside other applications, motivating
sub-10 MB model sizes after quantisation.

\subsection{Why Convolutional Neural Networks}

Deep learning, and convolutional neural networks (CNNs) in particular, have
transformed computer vision across many agricultural domains. Since the
landmark work of Mohanty \emph{et al.}~\cite{mohanty2016}, who demonstrated
99.35\% accuracy on the PlantVillage dataset using AlexNet and GoogLeNet,
researchers have progressively applied increasingly efficient architectures to
real-world, field-captured plant images \cite{alvarado2025}. Transfer
learning, which involves initialising a network with weights pretrained on
large-scale datasets such as ImageNet and fine-tuning on a domain-specific
collection, has emerged as the dominant strategy when labelled agricultural
data is limited \cite{transfer2024}. The advantage is twofold: low-level
features (edges, textures, colour gradients) generalise across natural images,
and the optimisation surface near a pre-trained minimum is dramatically more
benign than the one near random initialisation, allowing convergence in tens
of epochs with hundreds rather than millions of training images per class.

\subsection{Research Questions}\label{ssec:rq}

The deployment context outlined above translates into four concrete research
questions that this paper attempts to answer empirically rather than
rhetorically:

\begin{description}
\item[\textbf{RQ1}] On a single, imbalanced, field-captured cocoa-pod dataset,
how much accuracy and macro F1 does each architecture (ResNet-50,
EfficientNet-B0, MobileNetV3-Small) lose relative to the heaviest baseline
when all other variables — augmentation, optimiser, schedule, seed — are held
fixed?

\item[\textbf{RQ2}] Where does that loss concentrate? Is the cost of going
small spread evenly across classes, or is it absorbed disproportionately by
the rare classes that matter most for early disease detection?

\item[\textbf{RQ3}] Does the rank ordering of probability scores degrade as
fast as the argmax decision? In other words, can a calibrated top-$k$
interface partially recover the accuracy a smaller backbone gives up?

\item[\textbf{RQ4}] Under the deployment constraints of a 2--4\,GB Android
device with intermittent connectivity, which architecture is feasible after
INT8 quantisation, and what is the headroom for additional in-app stages
(detector, tracker, advisory layer)?
\end{description}

These questions structure the experimental setup in Sec.~\ref{sec:exp} and
the results in Sec.~\ref{sec:results}; each RQ is revisited explicitly in the
dedicated Discussion of Sec.~\ref{sec:discussion}.

\subsection{Contributions}

This paper makes the following contributions:

\begin{enumerate}
\item A complete deep-learning pipeline for multi-class cocoa pod disease
classification, converting bounding-box-annotated YOLO-format data into
cropped, class-balanced training samples.
\item A systematic comparative evaluation of three architectures
(ResNet-50~\cite{he2016}, EfficientNet-B0~\cite{tan2019}, and
MobileNetV3-Small~\cite{howard2019}) under \emph{identical} training
hyperparameters, augmentation, and class-weighting schemes, isolating
architecture as the only varying factor.
\item An extended evaluation that goes beyond aggregate accuracy to report
per-class precision, recall, F1, support, ROC AUC (one-vs-rest, micro, and
macro), and inference latency distribution (mean/median/P95/P99).
\item A model-footprint analysis covering FP32 size, estimated INT8 size, and
parameter count, mapping each architecture to a deployment tier (cloud
inference, on-device inference, ultra-low-end inference).
\item Practical recommendations for deploying the best edge-suitable model in
Ghanaian smallholder farming contexts, with explicit discussion of failure
modes and limitations.
\item A research-question-aligned discussion (Sec.~\ref{sec:discussion}) that
maps each empirical finding to a concrete deployment decision, rather than
ending the paper at an aggregate metrics table.
\end{enumerate}

To make the scope of these contributions concrete,
Table~\ref{tab:contributions} summarises which sections deliver each
contribution and which research question it answers; the table doubles as a
roadmap for a reader who wants to navigate the paper non-linearly.

\begin{table}[t]
\centering
\caption{Roadmap of contributions, anchoring sections, and research questions.}
\label{tab:contributions}
\renewcommand{\arraystretch}{1.15}
\begin{tabular}{p{0.04\columnwidth}p{0.45\columnwidth}p{0.18\columnwidth}p{0.12\columnwidth}}
\toprule
\textbf{\#} & \textbf{Contribution} & \textbf{Section} & \textbf{RQ}\\
\midrule
C1 & Pipeline (YOLO crops $\rightarrow$ class-balanced trainer) & \ref{sec:method} & --- \\
C2 & Identical-protocol architecture comparison & \ref{sec:exp},\ref{sec:results} & RQ1 \\
C3 & Per-class metrics + ROC + latency distribution & \ref{sec:results} & RQ2--3 \\
C4 & Footprint $\rightarrow$ deployment tier mapping & \ref{sec:results} & RQ4 \\
C5 & Smallholder-deployment recommendations & \ref{sec:discussion} & RQ4 \\
\bottomrule
\end{tabular}
\end{table}

The remainder of this paper is structured as follows.
Section~\ref{sec:related} reviews related work, including class-imbalance
strategies and the recent edge-deployment literature for agriculture.
Section~\ref{sec:background} formalises the multi-class classification
problem, the inductive biases of the chosen architectures, and the loss used.
Section~\ref{sec:method} describes the dataset, preprocessing pipeline, and
model architectures.
Section~\ref{sec:exp} details the experimental setup, including the
quantisation protocol used to derive INT8 footprints.
Section~\ref{sec:results} presents the per-architecture metric breakdown.
Section~\ref{sec:discussion} maps the results onto each of the four research
questions and onto the deployment recommendation for Ghanaian smallholders.
Section~\ref{sec:limits} reflects on limitations and threats to validity.
Section~\ref{sec:conclusion} concludes with directions for future work.
A multi-section appendix provides per-model detailed figures, full
reproducibility data, and per-class confusion counts.

%==============================================================
\section{Related Work}\label{sec:related}

\subsection{Plant Disease Detection with Deep Learning}

The systematic application of deep learning to plant disease detection began
in earnest with Mohanty \emph{et al.}~\cite{mohanty2016}, who trained CNN
classifiers on 54{,}306 images spanning 26 diseases across 14 crop species
from the PlantVillage dataset, achieving 99.35\% accuracy under controlled
laboratory conditions. Subsequent work consistently showed that performance
degrades significantly on field-captured images due to variations in
lighting, occlusion, background complexity, and disease stage, motivating
domain-specific data collection and augmentation \cite{alvarado2025}.
Two patterns dominate the post-2018 literature. First, the family of model
architectures has steadily shifted from heavy generic backbones (VGG-16,
GoogLeNet) toward families designed with efficiency in mind (EfficientNet,
MobileNet) or with stronger inductive biases (vision transformers).
Second, evaluation has moved from raw accuracy on lab images to
balanced-aware metrics on field images, reflecting the realisation that the
class imbalance present in nearly every agricultural dataset makes accuracy
a misleading single-number summary.

\subsection{Transfer Learning for Agricultural Imagery}

Transfer learning has become the standard approach for agricultural image
classification, where per-class sample counts are often insufficient to train
large networks from scratch \cite{transfer2024}. Pre-training on ImageNet
provides initialisation that captures universal low-level features such as
edges, textures, and colours that are transferable to plant imagery.
Fine-tuning the final classification layers, or all layers at a reduced
learning rate, consistently outperforms training from random initialisation
when datasets contain fewer than 10{,}000 samples per class \cite{anduyan2025}.
A complementary line of work investigates whether more domain-aligned
pre-training (for example, on iNaturalist or PlantCLEF) yields further gains;
the consensus is that ImageNet weights remain a strong default for crop-level
disease classification, with at most marginal additional benefit from
domain-aligned pre-training when the target dataset is small.

\subsection{Cocoa-Specific Disease Detection}

Research on automated cocoa disease detection has grown substantially since
2020. Buena\~{n}o Vera \emph{et al.}~\cite{buenano2024} evaluated
EfficientDet-Lite4 for simultaneous detection and localisation of monilia
(frosty pod rot) and black pod disease on cocoa pods, demonstrating that
lightweight detection models can achieve acceptable accuracy suitable for
mobile deployment. Their work targets a two-class detection problem rather
than a five-way classification and does not measure macro F1 over rare
classes. Anduyan \emph{et al.}~\cite{anduyan2025} conducted the most
comprehensive architecture comparison on cocoa pod images to date, evaluating
VGG16, VGG19, ResNet-50, ResNet-101, InceptionV3, and Xception as base
learners within three ensemble schemes; bagging achieved 100\% test accuracy
on a balanced 6{,}000-image dataset, though the study did not evaluate
lightweight architectures for edge deployment, and the balanced dataset does
not reflect the long-tailed reality of field collections.

Pagaduan \emph{et al.}~\cite{pagaduan2026} developed a mobile application for
cacao disease identification using deep learning on Filipino cocoa farms,
corroborating the challenge of high recall on rare disease classes.
Ayikpa \emph{et al.}~\cite{ayikpa2025} proposed a hybrid CNN and hand-crafted
feature approach, achieving 91.48\% accuracy; stacking further improved this
to 93.73\%, demonstrating that feature fusion can complement end-to-end
learning when data is scarce. Kumar \emph{et al.}~\cite{kumar2023} applied
DenseNet-19 to detect vascular streak dieback on cocoa leaf images, achieving
99.1\% accuracy on a 1{,}200-image dataset, although their leaf-only setting
sidesteps the geometric variation of pods that field deployment must handle.
Alvarado \emph{et al.}~\cite{alvarado2025} conducted a systematic literature
review of 86 studies on computer-vision methods for cocoa disease detection
from 2010 to 2025, confirming that CNN-based transfer learning dominates the
field and identifying a gap in lightweight, edge-deployable models.

\subsection{Lightweight Architectures for Edge Deployment}

The resource constraints of smartphone hardware motivate the use of compact
architectures. MobileNetV3-Small \cite{howard2019}, with its
depthwise-separable convolutions and hard-swish activations, achieves
competitive accuracy at 1.5 million parameters and sub-5\,ms inference on
mobile CPUs. EfficientNet-B0 \cite{tan2019} uses a compound scaling law that
simultaneously scales depth, width, and resolution. Both have been validated
for agricultural disease detection in recent benchmarks \cite{edgeai2024},
and are natural candidates for smartphone-based farm advisory systems.
ResNet-50 \cite{he2016} is included not because it is realistic for on-device
inference, but because it serves as a well-understood high-capacity baseline
that allows the cost of moving to a smaller architecture to be quantified.

\subsection{Class Imbalance Strategies in Agricultural Datasets}

Field-collected agricultural datasets are almost universally long-tailed: in
practice the majority of pods, leaves, or fruits photographed during a crop
campaign are healthy, and individual disease categories collectively occupy
only a fraction of the data. The literature offers four broadly distinct
strategies for handling this imbalance.

\textbf{Re-weighting.} Loss-level re-weighting multiplies each sample's loss
by an inverse-frequency-style weight so that minority classes contribute more
gradient per sample~\cite{cuiclassbalanced2019}. This is the strategy adopted
in the present paper. Its appeal is that it requires no change to the data
pipeline and no change to the model architecture; it modifies only the
training objective. The cost is that very rare classes can dominate the
gradient and produce noisy updates if their per-batch presence is also
volatile.

\textbf{Re-sampling.} Re-sampling either oversamples the minority class
(repeating samples or generating synthetic ones via SMOTE-style
interpolation~\cite{chawlasmote2002}) or undersamples the majority class. The
former increases epoch length and risks overfitting to repeated minority
samples; the latter discards information from the majority class.
Bagging-based ensembles such as the one Anduyan \emph{et al.}~\cite{anduyan2025}
deploy on cocoa images implicitly re-sample at the bootstrap level.

\textbf{Margin- and focal-style losses.} Focal loss~\cite{linfocal2017}
down-weights samples on which the model is already confident, focusing
gradient capacity on hard examples that often happen to be minority-class
samples. Class-balanced focal loss~\cite{cuiclassbalanced2019} extends this
idea with an effective-sample-count term that interpolates between uniform
weighting and inverse-frequency weighting. Margin-based losses such as LDAM
adjust the decision boundary itself per class.

\textbf{Two-stage pipelines.} Where a detector precedes a classifier, the
detector's confidence threshold can act as an implicit re-balancer by
filtering out trivially easy backgrounds before the classifier ever sees
them. Buena\~no Vera \emph{et al.}~\cite{buenano2024} use exactly this
pattern with EfficientDet-Lite4 in front of a binary classifier. The
trade-off is system complexity and the brittleness of compounding two error
sources.

This paper adopts inverse-frequency re-weighting for two reasons that the
related work does not jointly satisfy. First, it allows the architecture
comparison to be conducted under a single, lightweight protocol rather than
inside a detector--classifier sandwich whose detector behaviour would itself
become an uncontrolled variable. Second, it leaves room — explored in
Sec.~\ref{sec:limits} — for a follow-up paper to ablate focal and
class-balanced focal losses on the same split as a controlled comparison.

\subsection{Edge Inference and Model Compression for Agriculture}

A separate strand of recent work focuses on the engineering of running
neural networks on smartphone-class hardware in low-resource settings.
Howard \emph{et al.}~\cite{howard2019} introduced MobileNetV3 with two
sibling architectures, MobileNetV3-Large and MobileNetV3-Small, designed
explicitly for mobile inference; the small variant — the one used in this
paper — was target-tuned for sub-5\,ms inference on commodity ARM CPUs.
Tan and Le~\cite{tan2019} introduced EfficientNet, whose compound-scaling
law produces a family of architectures that trade accuracy for compute along
a controllable Pareto curve. Krizhevsky's seminal AlexNet~\cite{krizhevsky2012}
and Simonyan and Zisserman's VGG family~\cite{simonyan2015} continue to be
referenced as the heavy upper bound, but neither is plausible for offline
on-device inference in 2026.

For the deployment side specifically, three quantisation regimes dominate
the recent literature. Post-training INT8 quantisation, the simplest, takes a
trained FP32 network and converts both weights and activations to 8-bit
integers using a per-channel scale factor; for the architectures studied
here this nominally yields a $4\times$ memory reduction and, on hardware
with INT8 SIMD support, a meaningful latency reduction as well. Quantisation-aware
training (QAT) inserts simulated quantisation operators during fine-tuning,
recovering most of the accuracy lost by post-training quantisation but
requiring a second training pass. Mixed-precision quantisation
(weight-INT4 + activation-INT8, for example) trades a small additional
accuracy hit for further footprint reduction. This paper reports estimated
INT8 sizes on the basis of a $4\times$ post-training-quantisation factor and
flags QAT validation as future work; precise on-device validation requires a
target-hardware bench that is outside the scope of this study.

A complementary direction is structured pruning, in which entire convolution
filters are removed and the remaining weights are fine-tuned to recover
accuracy. For MobileNetV3-Small, channel pruning is known to interact poorly
with the depthwise-separable structure because the pointwise convolutions
already implement most of the channel mixing; literature reviews on
edge-AI for agriculture~\cite{edgeai2024} consistently report that
quantisation alone yields larger, more reliable footprint reductions than
pruning on this family of architectures.

\subsection{Knowledge Gap Addressed by This Work}

Three structural gaps emerge when the threads above are read together.
\textbf{First}, no published study to our knowledge runs a heavy baseline,
a Pareto-optimal mid-size architecture, and a mobile-first architecture
through \emph{the same training protocol on the same imbalanced cocoa-pod
dataset}, leaving the cost of each compression step systematically
unmeasured. \textbf{Second}, the existing cocoa-deep-learning literature
is dominated by aggregate accuracy on balanced subsets and rarely reports
per-class macro F1, ROC, or rare-class recall — the precise metrics that
matter for a farmer-facing system. \textbf{Third}, the deployment-feasibility
question (offline, mid-range Android, INT8) is typically addressed in
isolation, not as part of an architecture comparison, so a Ghanaian
deployment team has no apples-to-apples evidence on which to base a
choice. This paper attempts to close all three gaps simultaneously on a
single dataset and a single training protocol.

\subsection{Positioning of This Work}

Across these threads, the literature splits into two camps: balanced-dataset
benchmarks where heavy architectures hit very high accuracy, and
mobile-deployment papers that run a single lightweight architecture without
the heavy baseline alongside. This work spans both camps on \emph{the same}
imbalanced field dataset: it puts a heavy backbone, a balanced
parameter--accuracy compound-scaled backbone, and a mobile-first backbone
through the same protocol so that the cost of going small can be measured per
class rather than only in aggregate. The headline number is therefore not
``what is our best accuracy'' but ``what does each architecture lose, and
where, when we shrink it.''

%==============================================================
\section{Background and Problem Formulation}\label{sec:background}

\subsection{Multi-Class Classification Setup}

Let $\mathcal{X} \subset \mathbb{R}^{H\times W\times 3}$ be the space of RGB
pod crops (after the bounding-box extraction described in
Sec.~\ref{sec:method}) and let $\mathcal{Y}=\{0,1,2,3,4\}$ denote the five
target labels: Healthy, \textit{C. foraseminis}, \textit{M. perniciosa},
\textit{M. roreri}, and \textit{P. palmivora}. Given a labelled training set
$\mathcal{D}_{\text{train}}=\{(x_i,y_i)\}_{i=1}^{N}$, we learn a function
$f_\theta:\mathcal{X}\rightarrow\mathbb{R}^{|\mathcal{Y}|}$ that produces a
vector of class logits, which is mapped to a probability vector via softmax:
\begin{equation}
p_\theta(y=c\mid x)=\frac{\exp(f_\theta(x)_c)}
{\sum_{c'=0}^{4}\exp(f_\theta(x)_{c'})}.
\end{equation}
At test time the predicted label is $\hat{y}=\arg\max_c p_\theta(y=c\mid x)$.

\subsection{Class-Weighted Cross-Entropy Loss}

The dataset exhibits substantial class imbalance (Table~\ref{tab:dataset}):
the Healthy class is roughly 27$\times$ larger than \textit{M. perniciosa}.
Plain cross-entropy in this regime drives the network to over-predict the
majority class because doing so is the lowest-loss strategy on average. To
counteract this, we weight each class by the inverse of its frequency:
\begin{equation}
w_c \;=\; \frac{N}{C\cdot n_c},
\qquad
\sum_{c=0}^{C-1} w_c\,n_c \;=\; N,
\label{eq:weights}
\end{equation}
where $C=5$ is the number of classes, $n_c$ is the number of training samples
in class $c$, and $N$ is the total training set size. The training objective
is the weighted cross-entropy:
\begin{equation}
\mathcal{L}(\theta)=-\frac{1}{N}\sum_{i=1}^{N} w_{y_i}\log p_\theta(y_i\mid x_i).
\label{eq:loss}
\end{equation}
For our split this produces weights $w=[0.295,2.737,6.386,1.755,1.940]$, so a
misclassified \textit{M. perniciosa} sample contributes about 21$\times$ the
gradient of a misclassified Healthy sample.

\subsection{Choice of Macro F1 as the Headline Metric}

Aggregate accuracy on this dataset is dominated by the Healthy class. A
trivial classifier that predicts ``Healthy'' for every sample reaches roughly
65\% test accuracy without learning anything about disease at all. Macro F1,
defined as
\begin{equation}
\text{F1}_{\text{macro}}=\frac{1}{C}\sum_{c=0}^{C-1}
\frac{2\,\mathrm{prec}_c\,\mathrm{rec}_c}{\mathrm{prec}_c+\mathrm{rec}_c},
\end{equation}
weights every class equally regardless of frequency and therefore directly
penalises a model for ignoring rare classes. Throughout this paper accuracy
is reported alongside macro F1, but architecture decisions are made primarily
on macro F1.

\subsection{Receiver Operating Characteristic and Probability Calibration}

A second, complementary view of the model is provided by the Receiver
Operating Characteristic (ROC) curve. For a one-vs-rest treatment of class
$c$, every sample $(x_i,y_i)$ generates a binary outcome $\mathbb{1}[y_i=c]$
and a real-valued score $p_\theta(y=c\mid x_i)$. Sweeping a decision
threshold $\tau\in[0,1]$ across the score yields a curve whose
$y$-coordinate is the true positive rate at $\tau$ and whose $x$-coordinate
is the false positive rate; the area under this curve (AUC) summarises the
\emph{rank-ordering} quality of the score independently of the chosen
threshold. AUC is therefore decoupled from the argmax decision: a model
that misclassifies under a 0.5 threshold but assigns a higher probability
to the correct class than to any incorrect one will still have AUC near 1.
For an interface that exposes top-$k$ candidates with a confidence
indicator (the deployment direction recommended in
Sec.~\ref{sec:discussion}), AUC is the more honest predictor of user
experience than accuracy.

\subsection{Inductive Biases of the Three Architectures}\label{ssec:bias}

The three architectures evaluated in this paper are not interchangeable
``black-box CNNs''; they differ in inductive bias in ways that determine
where each one is likely to succeed and fail on the cocoa-pod task.

\textbf{Residual learning} (ResNet-50). A residual block computes
$y = \mathcal{F}(x; W) + x$, where the identity shortcut $x$ permits
gradients to flow back through the network without compounding through every
non-linearity. Empirically this enables stable training of networks on the
order of fifty layers deep, which in turn gives ResNet-50 a large effective
receptive field and a high-capacity representation space.
For the cocoa task, ResNet-50's depth is a benefit on the dominant Healthy
class — where every additional layer can refine subtle texture features —
but a liability on the small minority classes, because deep over-parameterised
networks tend to memorise rare examples rather than generalise from them. The
mild train--val gap visible in our ResNet-50 training curves
(Fig.~\ref{fig:appA_resnet_curve}) is an instance of this tendency.

\textbf{Compound-scaled MBConv blocks} (EfficientNet-B0). Tan and
Le~\cite{tan2019} observed that depth, width, and input resolution should be
scaled jointly rather than independently, and proposed a compound coefficient
that does so under a fixed compute budget. EfficientNet-B0 — the smallest
member of the family — pairs this scaling with mobile-inverted bottleneck
convolutions (MBConv) that include a squeeze-and-excitation
attention module per block. The squeeze-and-excitation gate gives
EfficientNet-B0 a channel-wise attention signal that is well suited to
agricultural imagery, where the disease-discriminative features (lesions,
darkening, fluffy white growth) are localised to particular channels of
particular layers. Of the three architectures it is therefore the most
parameter-efficient on rare-class recall, a prediction the
results in Sec.~\ref{sec:results} confirm empirically.

\textbf{Depthwise-separable convolutions and platform-aware NAS}
(MobileNetV3-Small). MobileNetV3-Small replaces every standard convolution
with a depthwise convolution (one filter per input channel) followed by a
pointwise $1\!\times\!1$ convolution that mixes the channels. This
factorisation drops the per-layer compute by approximately
$1/k^2 + 1/c$ relative to a standard $k\!\times\!k$ convolution with $c$
output channels, which is what makes the architecture mobile-feasible.
The hard-swish activation $h\!\text{-}\!\mathrm{swish}(x)=x\cdot
\mathrm{ReLU6}(x+3)/6$ is a piecewise-linear approximation of the standard
swish that is cheaper to evaluate on integer hardware. The resulting
architecture is approximately $15\times$ smaller than ResNet-50 in parameter
count but, by construction, has substantially less channel mixing capacity
per layer; we expect — and the results confirm — a notable accuracy gap on
the rarest classes, where channel-level subtlety matters most.

\textbf{Why these three.} The three architectures span the
parameter--accuracy frontier of widely-used image classification backbones:
ResNet-50 represents the heavy reference point, EfficientNet-B0 represents
the modern efficiency-accuracy compromise, and MobileNetV3-Small represents
the mobile-first lower bound. Comparing them under an identical protocol on
the same imbalanced cocoa dataset isolates architecture as the only variable
and therefore turns ``how much accuracy does shrinking the model cost?''
from a rhetorical question into a measurable one.

%==============================================================
\section{Methodology}\label{sec:method}

\subsection{Dataset}

The dataset consists of 3{,}113 raw images of cocoa pods captured under
diverse field conditions, annotated using Label Studio with YOLO-format
bounding boxes. Following bounding-box extraction, a total of 3{,}285
individual pod crops were obtained (some images contain multiple pods). The
five target classes and their distributions are summarised in
Table~\ref{tab:dataset}, and visualised in Fig.~\ref{fig:class_dist}.

\begin{table}[t]
\centering
\caption{Dataset class distribution (after bounding-box crop extraction; seed=42).}
\label{tab:dataset}
\renewcommand{\arraystretch}{1.15}
\begin{tabular}{clrrr}
\toprule
\textbf{ID} & \textbf{Class} & \textbf{Train} & \textbf{Val} & \textbf{Test} \\
\midrule
0 & Healthy                    & 1{,}560 & 318 & 324 \\
1 & \textit{C. foraseminis}    &   168   &  46 &  41 \\
2 & \textit{M. perniciosa}     &    72   &  13 &  12 \\
3 & \textit{M. roreri}         &   262   &  53 &  66 \\
4 & \textit{P. palmivora}      &   237   &  62 &  51 \\
\midrule
\textbf{Total} & & \textbf{2{,}299} & \textbf{492} & \textbf{494} \\
\bottomrule
\end{tabular}
\end{table}

\begin{figure}[t]
\centering
\includegraphics[width=\columnwidth]{class_distribution.png}
\caption{Per-split class counts. The Healthy class dominates every split; the
rarest class (\textit{M.\ perniciosa}) has only 72 training samples,
motivating inverse-frequency loss weighting (Eq.~\ref{eq:weights}).}
\label{fig:class_dist}
\end{figure}

The dataset exhibits substantial class imbalance: the Healthy class
constitutes approximately 63\% of all samples, while \textit{M. perniciosa}
accounts for only 2.3\%. Fig.~\ref{fig:dataset_samples} presents one
representative field image from each class, illustrating the visual diversity
of pod conditions and the challenging lighting and background variability
that the model must handle.

\begin{figure*}[t]
\centering
\begin{subfigure}[b]{0.18\textwidth}
\includegraphics[width=\linewidth]{dataset_class0_healthy.jpg}
\caption{Healthy}
\end{subfigure}\hfill
\begin{subfigure}[b]{0.18\textwidth}
\includegraphics[width=\linewidth]{dataset_class1_borer.jpg}
\caption{\textit{C. foraseminis}}
\end{subfigure}\hfill
\begin{subfigure}[b]{0.18\textwidth}
\includegraphics[width=\linewidth]{dataset_class2_witches.jpg}
\caption{\textit{M. perniciosa}}
\end{subfigure}\hfill
\begin{subfigure}[b]{0.18\textwidth}
\includegraphics[width=\linewidth]{dataset_class3_frosty.jpg}
\caption{\textit{M. roreri}}
\end{subfigure}\hfill
\begin{subfigure}[b]{0.18\textwidth}
\includegraphics[width=\linewidth]{dataset_class4_blackpod.jpg}
\caption{\textit{P. palmivora}}
\end{subfigure}
\caption{Representative field images from each of the five target classes.
Images are captured under diverse lighting and background conditions typical
of Ghanaian cocoa farms.}
\label{fig:dataset_samples}
\end{figure*}

\subsection{Data Preprocessing Pipeline}

\textbf{Bounding-box crop extraction.}
Each YOLO annotation provides the normalised centre coordinates $(x_c,y_c)$
and dimensions $(w,h)$ of a pod bounding box. Absolute pixel coordinates with
a 5\% padding factor $p$ are computed as
\begin{align}
x_1 &= \max\!\bigl(0,\;(x_c - w/2 - p)\cdot W\bigr),\\
y_1 &= \max\!\bigl(0,\;(y_c - h/2 - p)\cdot H\bigr),\\
x_2 &= \min\!\bigl(W,\;(x_c + w/2 + p)\cdot W\bigr),\\
y_2 &= \min\!\bigl(H,\;(y_c + h/2 + p)\cdot H\bigr),
\end{align}
where $W$ and $H$ are the image width and height. The pod crop $I[y_1{:}y_2,
x_1{:}x_2]$ is converted to a square 256$\times$256 image by anisotropic
resize (the same convention used by torchvision's reference recipes).

\textbf{Augmentation.} Training images are passed through random crop
($256\!\times\!256$ resize, $224\!\times\!224$ random crop), horizontal and
vertical flips, rotation up to $\pm15^\circ$, and colour jitter (brightness,
contrast, saturation $\pm0.3$; hue $\pm0.1$). Validation and test images are
resized directly to $224\!\times\!224$ with no augmentation.

\textbf{Normalisation.} All images are normalised using ImageNet channel
means $\boldsymbol{\mu}=[0.485,0.456,0.406]$ and standard deviations
$\boldsymbol{\sigma}=[0.229,0.224,0.225]$ to match the input distribution
expected by the pre-trained backbones. Table~\ref{tab:augparams} summarises
the full set of augmentation parameters.

\begin{table}[t]
\centering
\caption{Data augmentation parameters used for training.}
\label{tab:augparams}
\renewcommand{\arraystretch}{1.15}
\begin{tabular}{ll}
\toprule
\textbf{Operation} & \textbf{Setting} \\
\midrule
Resize                 & $256\times256$ \\
Random crop            & $224\times224$ \\
Horizontal flip        & probability 0.5 \\
Vertical flip          & probability 0.5 \\
Rotation               & $\pm15^\circ$ \\
Colour jitter (b/c/s)  & $\pm0.3$ \\
Colour jitter (hue)    & $\pm0.1$ \\
Channel mean ($\mu$)   & $[0.485,0.456,0.406]$ \\
Channel std ($\sigma$) & $[0.229,0.224,0.225]$ \\
\bottomrule
\end{tabular}
\end{table}

\subsection{Model Architectures}

All three models are initialised with ImageNet pre-trained weights and
adapted for five-class classification by replacing their final layers. Where
the original backbone defines a 1000-way classification head, that head is
replaced with a freshly initialised five-way linear layer; all other
parameters are fine-tuned from their pre-trained values.

\textbf{ResNet-50}~\cite{he2016} uses residual skip connections enabling
effective training of deep networks. It consists of four residual stages
followed by global average pooling and the (replaced) five-way linear output.
ResNet-50 contains 23.52 million trainable parameters in our configuration.

\textbf{EfficientNet-B0}~\cite{tan2019} uses mobile-inverted bottleneck
convolution (MBConv) blocks with squeeze-and-excitation modules, scaled
uniformly via a compound coefficient. After replacing the final classifier, it
has 4.01 million trainable parameters.

\textbf{MobileNetV3-Small}~\cite{howard2019} employs depthwise-separable
convolutions, hard-swish activations, and platform-aware neural architecture
search for mobile deployment. It has 1.52 million trainable parameters,
approximately 15$\times$ fewer than ResNet-50, and is the only architecture
in this comparison that fits comfortably alongside a typical Android app
budget after INT8 quantisation.

Table~\ref{tab:arch} summarises the headline architecture statistics.

\begin{table}[t]
\centering
\caption{Architecture summary. FP32 sizes are measured directly from the
state dict; INT8 estimates assume 4$\times$ post-training quantisation
without weight pruning.}
\label{tab:arch}
\renewcommand{\arraystretch}{1.15}
\begin{tabular}{lrrrr}
\toprule
\textbf{Model} & \textbf{Params (M)} & \textbf{FP32 (MB)} & \textbf{INT8 (MB)} & \textbf{Depth} \\
\midrule
ResNet-50          & 23.52 & 89.7 & 22.4 & 50 \\
EfficientNet-B0    &  4.01 & 15.3 &  3.8 & $\sim$237 \\
MobileNetV3-Small  &  1.52 &  5.8 &  1.5 & $\sim$73 \\
\bottomrule
\end{tabular}
\end{table}

\subsection{Why Inverse-Frequency Weighting}\label{ssec:invfreq}

The choice of inverse-frequency cross-entropy in Eq.~\ref{eq:loss} is
deliberate but not arbitrary. Let $\pi_c = n_c/N$ denote the empirical class
prior on the training set. Plain (unweighted) cross-entropy minimises an
estimate of the expected loss under $\pi$, which on this dataset means an
estimator that is heavily biased toward predicting the dominant class.
Multiplying each loss term by $w_c = 1/(C\pi_c)$ rescales the per-class
contribution to be uniform, which is equivalent — up to a constant — to
minimising the expected loss under a uniform class prior $\pi'_c = 1/C$.
This is the implicit assumption a macro-F1-driven evaluation makes.

The downside of this rescaling is that it inflates the gradient variance for
rare classes. With $w_{\text{witches}}\!=\!6.39$ and only 72 training
samples, the contribution of any one mis-labelled witches sample to a
mini-batch gradient is roughly $21\times$ that of a healthy sample. We
mitigate this by using Adam's adaptive per-parameter learning rate, which
absorbs gradient-magnitude variation across mini-batches without further
tuning. An ablation against focal loss and class-balanced focal loss
remains a follow-up worth the effort but is outside the scope of this paper.

\subsection{Alternative Class-Balancing Strategies Considered}\label{ssec:alt}

Three alternatives to inverse-frequency weighting were considered before
finalising the protocol:

\textbf{Oversampling.} Repeating minority-class samples up to majority size
inflates epoch length by roughly $4\times$ on this dataset and risks
overfitting to the duplicated rare samples — an outcome that is
particularly harmful when those samples are noisy in the first place.

\textbf{Synthetic minority oversampling (SMOTE).} SMOTE interpolates new
samples between existing minority-class points in feature space.
SMOTE was designed for tabular data and translates poorly to natural
images, where straight-line interpolation between two cocoa pods does not
produce a third realistic cocoa pod.

\textbf{Focal loss.} Focal loss~\cite{linfocal2017} down-weights easy
samples by $(1-p_t)^\gamma$. On a 72-sample minority class, focal loss and
inverse-frequency weighting are nearly equivalent in expectation, but focal
loss adds a hyperparameter ($\gamma$) that must itself be tuned. We chose
the simpler scheme to keep architecture comparison clean.

\subsection{Training Strategy}

All models are trained under identical hyperparameters
(Table~\ref{tab:hyperparams}) using the Adam optimiser with weight decay
$\lambda=10^{-4}$ and a cosine-annealing learning-rate schedule from
$\eta_0=10^{-3}$. Training terminates early when validation accuracy fails to
improve for seven consecutive epochs. The same random seed (42) governs the
PyTorch CPU/GPU RNGs and the dataset split, so the train/val/test partition
is identical across all three runs.

\begin{table}[t]
\centering
\caption{Training hyperparameters (identical across all three architectures).}
\label{tab:hyperparams}
\renewcommand{\arraystretch}{1.15}
\begin{tabular}{ll}
\toprule
\textbf{Hyperparameter} & \textbf{Value} \\
\midrule
Optimiser              & Adam \\
Initial learning rate  & $1\times10^{-3}$ \\
Weight decay           & $1\times10^{-4}$ \\
LR schedule            & Cosine annealing \\
Batch size             & 32 \\
Max epochs             & 40 \\
Early stopping patience & 7 \\
Input resolution       & $224\times224$ \\
Loss function          & Weighted cross-entropy \\
Pretrained weights     & ImageNet-1K \\
Random seed            & 42 \\
\bottomrule
\end{tabular}
\end{table}

%==============================================================
\section{Experimental Setup}\label{sec:exp}

\subsection{Implementation and Hardware}

All experiments are implemented in Python 3.9 using PyTorch 2.8 and
torchvision. Training uses Apple Metal Performance Shaders (MPS) GPU
acceleration on an Apple M-series processor with 16~GB unified memory. The
choice of Apple Silicon as the development platform is deliberate: it
mirrors, in throughput class, the modern mid-range Android devices on which
the model is ultimately intended to run, so that latency numbers reported
here are within an order of magnitude of the on-device deployment target.

\subsection{Inference Latency Protocol}

Per-image inference latency is measured with a 1-image batch over 200 timed
forward passes. Each timing call is bracketed by hardware synchronisation
(\texttt{torch.mps.synchronize}) so that the asynchronous launch of MPS
kernels does not flatter the measurement. A 20-pass warm-up precedes timing.
From these 200 samples we report mean, median, standard deviation, P95, and
P99. The latency \emph{distribution} matters more than the mean for a real
deployment: a model whose 99th-percentile latency is 5$\times$ the mean will
visibly stutter even if its average looks acceptable.

\subsection{Evaluation Metrics}

Models are evaluated on the held-out test set ($N=494$) using the following
metrics:
\begin{itemize}
\item \textbf{Overall accuracy}: $\sum_i\mathbb{1}[\hat{y}_i=y_i]/N$.
\item \textbf{Macro F1}: unweighted mean of per-class F1 (Eq.~3).
\item \textbf{Macro precision and recall}.
\item \textbf{Per-class precision/recall/F1/support}.
\item \textbf{Confusion matrix} (raw counts and row-normalised).
\item \textbf{ROC AUC} per class (one-vs-rest), micro-averaged, and
macro-averaged.
\item \textbf{Inference latency} (ms per image): mean, median, P95, P99.
\item \textbf{Parameter count} and FP32/INT8 model size.
\end{itemize}
Macro-averaged metrics are prioritised over accuracy because the test set is
class-imbalanced (Healthy comprises 65.6\% of test samples).

\subsection{Dataset Split}

The full set of 3{,}285 annotated crops is divided into training (70\%),
validation (15\%), and test (15\%) sets using stratified random sampling with
seed 42. The exact per-class counts are reported in Table~\ref{tab:dataset}
and visualised in Fig.~\ref{fig:class_dist}; the same split is used for all
three architectures.

\subsection{Quantisation Methodology}\label{ssec:quant}

The INT8 footprints reported throughout this paper are derived from a single
methodology that deserves explicit description because it determines whether
the headline ``edge-feasible'' claim survives scrutiny. Three points apply.

\textbf{Per-channel symmetric weight quantisation.} Each weight tensor is
quantised channel-wise with a symmetric scale $s_c$ chosen so that the
weight range $[\min_c, \max_c]$ maps to $[-127, 127]$. Per-channel
symmetric quantisation is the default in PyTorch's
\texttt{torch.quantization} backend and in TFLite's INT8 export path; it
produces no accuracy degradation beyond what fully-symmetric INT8 imposes
on the activation path.

\textbf{Activation calibration.} Activations are quantised with per-tensor
asymmetric scales calibrated on a 200-image subset of the validation set.
Asymmetric quantisation is appropriate because depthwise convolutions and
hard-swish activations produce activation distributions that are visibly
non-zero-centred.

\textbf{Reported INT8 sizes are estimates.} The footprint values in
Table~\ref{tab:arch} and Table~\ref{tab:deploy} are computed as
$0.25\times$ the FP32 size; this is an upper bound on the actual quantised
weight footprint because batch-norm parameters and biases are typically kept
in FP32 even after weight quantisation. A bench on actual target hardware
(see Sec.~\ref{sec:limits}) is required before the footprint claims are
elevated from estimate to deployment SLO.

\subsection{Reproducibility Engineering}

A reproducible result is one in which an independent reader, given the same
data and the same code, observes identical metrics. Achieving this on a
modern deep-learning stack requires more discipline than seeding a single
RNG. The protocol used here pins:
\begin{itemize}
\item the dataset split (stratified shuffle with seed 42, recomputed
deterministically each run);
\item the PyTorch CPU and GPU RNGs (seeded via
\texttt{torch.manual\_seed} and \texttt{torch.mps.manual\_seed});
\item the NumPy RNG used by all augmentation transforms;
\item the cuDNN/MPS deterministic flag (where available; MPS support is
partial as of PyTorch 2.8 and may introduce sub-percent-point variance
across machines);
\item the Python and library versions enumerated in
Table~\ref{tab:repro}.
\end{itemize}
This protocol is sufficient for byte-identical metrics on the development
machine across reruns, and within ${\sim}0.3$ accuracy points across
machines with different MPS hardware revisions. The exact code, including
the deterministic-RNG configuration block, is reproduced in the appendix
listing.

\subsection{System Pipeline}

Fig.~\ref{fig:pipeline} summarises the end-to-end pipeline. Raw smartphone
images annotated with YOLO-format bounding boxes are processed by the crop
extractor to produce per-pod crops, which are split deterministically into
train/val/test, augmented (training only) and normalised, then fed into the
selected backbone. The backbone produces five-way logits, which are converted
to probabilities by softmax and to a discrete prediction by argmax. At
training time the loss is the inverse-frequency-weighted cross-entropy of
Eq.~\ref{eq:loss}; at inference time only the forward path runs, with the
augmentation block bypassed.

\begin{figure*}[t]
\centering
\includegraphics[width=0.95\textwidth]{fig3_pipeline.png}
\caption{End-to-end inference and training pipeline. Stage~1 (input pipeline,
CPU) extracts pod crops from YOLO bounding boxes, resizes to
$224\!\times\!224$, and applies augmentation; Stage~2 (classifier) normalises
with ImageNet statistics, runs the chosen backbone, and emits a five-way
softmax decision. Augmentation is applied only on training samples; validation
and test samples are resized directly to $224\!\times\!224$. The training
objective is the inverse-frequency-weighted cross-entropy of
Eq.~\ref{eq:loss}.}
\label{fig:pipeline}
\end{figure*}

%==============================================================
\section{Results and Discussion}\label{sec:results}

\subsection{Overall Performance Comparison}

Table~\ref{tab:results} presents the test-set performance of all three
models, selected by maximum validation accuracy.

\begin{table*}[t]
\centering
\caption{Test-set performance comparison. Macro precision/recall/F1 are
unweighted means over the five classes; AUC is the macro one-vs-rest area
under the ROC curve.}
\label{tab:results}
\renewcommand{\arraystretch}{1.15}
\begin{tabular}{lccccccc}
\toprule
\textbf{Model} & \textbf{Acc.} & \textbf{Macro F1} & \textbf{Macro Prec.} & \textbf{Macro Rec.} & \textbf{AUC (macro)} & \textbf{Latency (ms)} & \textbf{Params (M)} \\
\midrule
ResNet-50          & 0.8725 & 0.7410 & 0.7285 & 0.7718 & 0.945 & 5.93 & 23.52 \\
EfficientNet-B0    & 0.8684 & \textbf{0.7547} & \textbf{0.7397} & \textbf{0.7774} & 0.939 & 5.41 & 4.01 \\
MobileNetV3-Small  & 0.8178 & 0.6745 & 0.6697 & 0.7413 & 0.929 & \textbf{3.88} & \textbf{1.52} \\
\bottomrule
\end{tabular}
\end{table*}

Fig.~\ref{fig:model_comparison} provides a visual summary of the performance
comparison. EfficientNet-B0 leads on macro F1 despite having 6$\times$ fewer
parameters than ResNet-50, confirming that compound scaling provides an
efficient accuracy--parameter trade-off for this task. The accuracy gap
between ResNet-50 and EfficientNet-B0 is only 0.41 percentage points, while
the parameter ratio is roughly 5.9$\times$, illustrating the diminishing
returns of additional depth on this dataset.

\begin{figure}[t]
\centering
\includegraphics[width=\columnwidth]{model_comparison.png}
\caption{Side-by-side comparison of all three architectures across accuracy,
macro F1, precision, and recall.}
\label{fig:model_comparison}
\end{figure}

\subsection{Per-Class Analysis}

The per-class breakdown (Table~\ref{tab:perclass}) reveals that the Healthy
class consistently achieves high F1 across all architectures.
\textit{M. perniciosa}, with only 72 training samples, remains the primary
bottleneck for macro F1 across all models. Notably, EfficientNet-B0 retains
0.643 F1 on this rare class, while MobileNetV3-Small collapses to 0.367,
losing $\sim$28 absolute points. This is consistent with the general finding
that smaller capacity networks pay for their efficiency disproportionately on
the rarest classes.

\begin{table*}[t]
\centering
\caption{Per-class F1 by model on the held-out test set. Bold marks the best
F1 in each row.}
\label{tab:perclass}
\renewcommand{\arraystretch}{1.15}
\begin{tabular}{lcccc}
\toprule
\textbf{Class} & \textbf{Support} & \textbf{ResNet-50} & \textbf{EfficientNet-B0} & \textbf{MobileNetV3-Small} \\
\midrule
Healthy                  & 324 & \textbf{0.9482} & 0.9406 & 0.9056 \\
\textit{C. foraseminis}  &  41 & 0.6579 & \textbf{0.6753} & 0.6053 \\
\textit{M. perniciosa}   &  12 & 0.5333 & \textbf{0.6429} & 0.3673 \\
\textit{M. roreri}       &  66 & 0.8320 & \textbf{0.8358} & 0.8333 \\
\textit{P. palmivora}    &  51 & \textbf{0.7333} & 0.6789 & 0.6610 \\
\midrule
\textbf{Macro avg}       & 494 & 0.7410 & \textbf{0.7547} & 0.6745 \\
\bottomrule
\end{tabular}
\end{table*}

Tables~\ref{tab:perclass_prec}~and~\ref{tab:perclass_rec} decompose F1 into
its precision and recall components. Two structural patterns emerge.
\textbf{First}, all three models show high precision and recall for Healthy,
which is unsurprising given its dominance and visual distinctiveness.
\textbf{Second}, on \textit{M. perniciosa}, recall is consistently higher
than precision: when the model is willing to commit to this rare class it is
often right (high recall), but it commits less often than the actual class
support warrants (low precision). This is a recoverable failure mode in
deployment because farmers can be shown a confidence score and a list of the
top-$k$ candidates rather than a hard top-1.

\begin{table}[t]
\centering
\caption{Per-class precision (test set).}
\label{tab:perclass_prec}
\renewcommand{\arraystretch}{1.15}
\begin{tabular}{lccc}
\toprule
\textbf{Class} & \textbf{ResNet-50} & \textbf{EfficientNet-B0} & \textbf{MobNetV3-S} \\
\midrule
Healthy           & 0.965 & 0.953 & 0.940 \\
Borer             & 0.714 & 0.722 & 0.657 \\
Witches           & 0.444 & 0.563 & 0.243 \\
Frosty            & 0.881 & 0.824 & 0.926 \\
Black pod         & 0.638 & 0.638 & 0.582 \\
\midrule
Macro             & 0.729 & 0.740 & 0.670 \\
\bottomrule
\end{tabular}
\end{table}

\begin{table}[t]
\centering
\caption{Per-class recall (test set).}
\label{tab:perclass_rec}
\renewcommand{\arraystretch}{1.15}
\begin{tabular}{lccc}
\toprule
\textbf{Class} & \textbf{ResNet-50} & \textbf{EfficientNet-B0} & \textbf{MobNetV3-S} \\
\midrule
Healthy           & 0.932 & 0.929 & 0.873 \\
Borer             & 0.610 & 0.634 & 0.561 \\
Witches           & 0.667 & 0.750 & 0.750 \\
Frosty            & 0.788 & 0.848 & 0.758 \\
Black pod         & 0.863 & 0.725 & 0.765 \\
\midrule
Macro             & 0.772 & 0.777 & 0.741 \\
\bottomrule
\end{tabular}
\end{table}

Fig.~\ref{fig:perclass_f1} and Fig.~\ref{fig:perclass_pr} visualise the
per-class metrics across all three models, making the rare-class gap
between EfficientNet-B0 and MobileNetV3-Small immediately obvious.

\begin{figure}[t]
\centering
\includegraphics[width=\columnwidth]{perclass_f1.png}
\caption{Per-class F1 across the three architectures. The gap between
EfficientNet-B0 and MobileNetV3-Small is largest on the rarest class
(\textit{M. perniciosa}, 12 test samples).}
\label{fig:perclass_f1}
\end{figure}

\begin{figure}[t]
\centering
\includegraphics[width=\columnwidth]{perclass_precrec.png}
\caption{Per-class precision (left) and recall (right) across the three
architectures.}
\label{fig:perclass_pr}
\end{figure}

\subsection{ROC Analysis}

Table~\ref{tab:auc} reports per-class one-vs-rest AUC for each architecture.
ResNet-50 and EfficientNet-B0 are essentially tied on macro AUC (0.945 vs
0.939), with MobileNetV3-Small slightly behind at 0.929. Importantly, all
three models clear AUC$\,>\,$0.86 on every class, including the rare
\textit{M.\ perniciosa}, indicating that the rank-ordering of probability is
informative even when the argmax decision is wrong; this matters for any
downstream calibration or thresholding stage.

\begin{table}[t]
\centering
\caption{Per-class ROC AUC (one-vs-rest). The macro-average appears in
Table~\ref{tab:results}.}
\label{tab:auc}
\renewcommand{\arraystretch}{1.15}
\begin{tabular}{lccc}
\toprule
\textbf{Class} & \textbf{ResNet-50} & \textbf{EffNet-B0} & \textbf{MobNetV3-S} \\
\midrule
Healthy           & 0.969 & 0.963 & 0.948 \\
\textit{C.\ foraseminis} & 0.925 & 0.897 & 0.868 \\
\textit{M.\ perniciosa} & 0.889 & 0.911 & 0.908 \\
\textit{M.\ roreri}     & 0.977 & 0.983 & 0.984 \\
\textit{P.\ palmivora}  & 0.964 & 0.943 & 0.936 \\
\midrule
Micro-avg               & 0.973 & 0.974 & 0.963 \\
\bottomrule
\end{tabular}
\end{table}

The per-architecture ROC curves are shown in
Figs.~\ref{fig:roc_resnet}--\ref{fig:roc_mobile}. They confirm that the
classification boundaries in probability space are crisp for the dominant
disease classes and somewhat softer for the rare ones, consistent with the
support-weighted analysis above.

\begin{figure}[t]
\centering
\includegraphics[width=0.95\columnwidth]{roc_ResNet50.png}
\caption{One-vs-rest ROC curves for ResNet-50 on the test set.}
\label{fig:roc_resnet}
\end{figure}

\begin{figure}[t]
\centering
\includegraphics[width=0.95\columnwidth]{roc_EfficientNet-B0.png}
\caption{One-vs-rest ROC curves for EfficientNet-B0 on the test set.}
\label{fig:roc_eff}
\end{figure}

\begin{figure}[t]
\centering
\includegraphics[width=0.95\columnwidth]{roc_MobileNetV3-Small.png}
\caption{One-vs-rest ROC curves for MobileNetV3-Small on the test set.}
\label{fig:roc_mobile}
\end{figure}

\subsection{Confusion Matrix Analysis}

The confusion matrices in Fig.~\ref{fig:cm_grid} reveal that
misclassifications occur primarily between visually similar disease classes.
\textit{P.\ palmivora} (black pod) and \textit{M.\ roreri} (frosty pod rot)
are the most commonly confused pair, as both manifest as external darkening
and lesion formation on the pod surface. \textit{C.\ foraseminis} (pod borer)
shows some confusion with healthy pods in early infestation stages, where
external lesions are subtle. The MobileNetV3-Small confusion matrix (rightmost
panel) shows a notable spillover from Healthy into the \textit{Witches}
class, accounting for its lower precision on that class --- an artifact of
the smaller network being unable to discriminate the more subtle morphological
cues.

\begin{figure*}[t]
\centering
\begin{subfigure}[b]{0.32\textwidth}
\includegraphics[width=\linewidth]{cm_norm_ResNet50.png}
\caption{ResNet-50 (normalised)}
\end{subfigure}\hfill
\begin{subfigure}[b]{0.32\textwidth}
\includegraphics[width=\linewidth]{cm_norm_EfficientNet-B0.png}
\caption{EfficientNet-B0 (normalised)}
\end{subfigure}\hfill
\begin{subfigure}[b]{0.32\textwidth}
\includegraphics[width=\linewidth]{cm_norm_MobileNetV3-Small.png}
\caption{MobileNetV3-Small (normalised)}
\end{subfigure}
\caption{Row-normalised confusion matrices for each architecture on the test
set. Diagonal cells are per-class recall; off-diagonals are confusion rates.}
\label{fig:cm_grid}
\end{figure*}

\subsection{Efficiency vs.\ Accuracy Trade-off}

The three architectures present a clear trade-off along the
accuracy--efficiency frontier (Fig.~\ref{fig:pareto}). ResNet-50 achieves the
highest classification performance at the cost of 23.52 million parameters,
making it impractical for offline smartphone inference. EfficientNet-B0
offers strong accuracy at 4 million parameters --- a 5.9$\times$ parameter
reduction vs.\ ResNet-50 with only a 0.41-point accuracy drop and a 1.4-point
\emph{gain} on macro F1. MobileNetV3-Small sacrifices a modest amount of
accuracy for a dramatic reduction in parameter count (1.52M) and inference
time (under 4\,ms), making it the most viable candidate for on-device
Ghanaian-farmer applications \cite{ayikpa2025,pagaduan2026}.
Fig.~\ref{fig:efficiency} illustrates this efficiency frontier explicitly,
plotting accuracy and macro F1 against model parameter count and inference
latency.

\begin{figure}[t]
\centering
\includegraphics[width=\columnwidth]{efficiency_comparison.png}
\caption{Efficiency vs.\ accuracy trade-off across the three architectures.
MobileNetV3-Small occupies a unique operating point combining sub-4\,ms
inference and 1.52M parameters with competitive accuracy.}
\label{fig:efficiency}
\end{figure}

\begin{figure}[t]
\centering
\includegraphics[width=\columnwidth]{pareto_f1_vs_params.png}
\caption{Macro F1 vs.\ parameter count on a log-x axis. Marker size is
proportional to parameter count. EfficientNet-B0 is Pareto-dominant: every
other point requires either more parameters for less F1 (ResNet-50) or
gives up more F1 for less compute (MobileNetV3-Small).}
\label{fig:pareto}
\end{figure}

\subsection{Inference Latency Distribution}

Fig.~\ref{fig:latency_dist} shows the per-image latency distribution over 200
runs for each architecture. Two practical observations matter for deployment.
First, MobileNetV3-Small's median is 3.6\,ms with P99 at 5.0\,ms --- a
worst-case better than EfficientNet-B0's mean. Second, the heavier networks
have tighter distributions, but their absolute floors are higher; for a
real-time scenario where each frame must complete inside 16.7\,ms (60\,fps),
all three are comfortably inside budget on Apple MPS, but the headroom is
4$\times$ larger for MobileNetV3-Small, which leaves room for a tracking or
cropping stage in front of the classifier. Table~\ref{tab:latency} reports
the full latency statistics including P95 and P99.

\begin{figure}[t]
\centering
\includegraphics[width=\columnwidth]{latency_distribution.png}
\caption{Box plot of per-image inference latency over 200 timed runs on Apple
MPS. The diamond marks the mean.}
\label{fig:latency_dist}
\end{figure}

\begin{table}[t]
\centering
\caption{Inference latency statistics (milliseconds) over 200 timed runs.
P95 and P99 are the empirical 95th and 99th percentiles.}
\label{tab:latency}
\renewcommand{\arraystretch}{1.15}
\begin{tabular}{lccccc}
\toprule
\textbf{Model} & \textbf{Mean} & \textbf{Median} & \textbf{Std} & \textbf{P95} & \textbf{P99} \\
\midrule
ResNet-50          & 5.93 & 5.92 & 0.04 & 6.02 & 6.10 \\
EfficientNet-B0    & 5.41 & 5.35 & 0.28 & 5.84 & 6.34 \\
MobileNetV3-Small  & 3.88 & 3.64 & 0.54 & 4.92 & 5.04 \\
\bottomrule
\end{tabular}
\end{table}

\subsection{Training Dynamics}

Training curves (Fig.~\ref{fig:curves}) show that all three models converge
within their allotted budget. EfficientNet-B0 exhibits the smoothest learning
curve with the least oscillation. ResNet-50 achieves the highest training
accuracy but shows a larger train--val gap, indicating mild overfitting that
early stopping partially mitigates. MobileNetV3-Small converges fastest and
triggers early stopping after only 10 epochs --- consistent with its smaller
capacity rapidly saturating on the available training data. The combined
training/validation grid for all three architectures appears in
Fig.~\ref{fig:curves} just before the References; per-model training
curves are reproduced in greater detail in Appendix~\ref{app:perm}.

\subsection{Comparison With Prior Cocoa-Disease Studies}

Table~\ref{tab:comparison} situates this work against recent cocoa-disease
publications. Headline accuracy on cocoa datasets ranges from $\sim$87\%
(this work, EfficientNet-B0 baseline of comparable parameter count) up to
100\% (Anduyan \emph{et al.}~\cite{anduyan2025}, balanced 6{,}000-image
dataset with bagging ensembles). Direct accuracy comparison across these rows
is unfair because the underlying datasets have very different class
distributions, image counts, and capture conditions; the table's purpose is
not to rank but to make the trade-offs visible.

\begin{table*}[t]
\centering
\caption{Comparison with recent cocoa-disease studies. Datasets and class
counts vary, so the ``Acc.'' column is not directly comparable across rows.}
\label{tab:comparison}
\renewcommand{\arraystretch}{1.15}
\begin{tabular}{p{0.32\textwidth}p{0.18\textwidth}rrp{0.18\textwidth}}
\toprule
\textbf{Study} & \textbf{Method} & \textbf{Cls.} & \textbf{Acc.} & \textbf{Edge focus} \\
\midrule
Buena\~no Vera \emph{et al.}~\cite{buenano2024} & EfficientDet-Lite4 (det.) & 2 & 0.78--0.85 & Yes \\
Anduyan \emph{et al.}~\cite{anduyan2025}        & VGG/ResNet/Inception ensemble (bagging) & 4 & 1.00 & No \\
Pagaduan \emph{et al.}~\cite{pagaduan2026}      & CNN (mobile)              & 4 & 0.92 & Yes \\
Ayikpa \emph{et al.}~\cite{ayikpa2025}          & CNN + handcrafted + stacking & 4 & 0.94 & No \\
Kumar \emph{et al.}~\cite{kumar2023}            & DenseNet-19 (leaf)        & 2 & 0.99 & No \\
\textbf{This work}                              & ResNet-50 / EfficientNet-B0 / MobileNetV3-Small & 5 & 0.82--0.87 & Yes \\
\bottomrule
\end{tabular}
\end{table*}

\subsection{Deployment Feasibility for Ghanaian Farmers}

The practical deployment scenario targets smallholder cocoa farmers in Ghana
who typically own mid-range Android smartphones with limited or intermittent
internet connectivity. Three constraints dominate model selection:
(i)~offline capability, (ii)~latency under 100\,ms, and
(iii)~memory footprint within a 2--4\,GB smartphone RAM budget.
MobileNetV3-Small satisfies all three constraints. At 1.52 million parameters,
the model occupies approximately 5.8\,MB in FP32 (or $\sim$1.5\,MB after 8-bit
quantisation). Its inference latency of under 4\,ms permits real-time disease
screening during field walks. PyTorch Mobile and ONNX Runtime both support
MobileNetV3-Small export, enabling straightforward Android and iOS
deployment. Table~\ref{tab:deploy} summarises the deployment classification
applied to each architecture.

\begin{table}[t]
\centering
\caption{Deployment tier for each architecture under the Ghanaian smallholder
constraints. ``Edge-OK'' = $<$10MB INT8 and $<$10ms median latency.}
\label{tab:deploy}
\renewcommand{\arraystretch}{1.15}
\begin{tabular}{lccc}
\toprule
\textbf{Model} & \textbf{INT8 (MB)} & \textbf{Median (ms)} & \textbf{Tier} \\
\midrule
ResNet-50          & 22.4 & 5.9 & Server / Cloud \\
EfficientNet-B0    &  3.8 & 5.4 & Edge-OK \\
MobileNetV3-Small  &  1.5 & 3.6 & Edge-Preferred \\
\bottomrule
\end{tabular}
\end{table}

\subsection{Training-Time Cost Analysis}

Beyond inference, the resources required to \emph{train} each architecture
also vary substantially. Table~\ref{tab:traintime} reports the wall-clock
training time and the number of optimizer steps before each model satisfied
the early-stopping criterion. Two findings are worth highlighting. First,
MobileNetV3-Small reaches its plateau after only 10 epochs ($\approx 720$
optimizer steps), roughly one-third of the training cost of ResNet-50.
Second, the per-epoch wall time on Apple MPS scales roughly linearly with
parameter count, but EfficientNet-B0 is closer in epoch time to ResNet-50
than its parameter count would suggest, because compound scaling produces a
deeper computation graph (more layers) even at smaller width. For research
groups iterating quickly under a fixed compute budget, this means that
EfficientNet-B0 is not necessarily a free lunch over ResNet-50 in
\emph{training} time, even though it dominates on inference and storage.

\begin{table}[t]
\centering
\caption{Approximate wall-clock training cost on Apple MPS (16 GB unified
memory). ``Steps'' is the cumulative number of optimizer updates.}
\label{tab:traintime}
\renewcommand{\arraystretch}{1.15}
\begin{tabular}{lccc}
\toprule
\textbf{Model} & \textbf{Epochs} & \textbf{Steps} & \textbf{Wall time} \\
\midrule
ResNet-50          & 37 & 2{,}664 & $\sim$41 min \\
EfficientNet-B0    & 18 & 1{,}296 & $\sim$22 min \\
MobileNetV3-Small  & 10 &   720   & $\sim$11 min \\
\bottomrule
\end{tabular}
\end{table}

\subsection{Field-Workflow Considerations}

A model good enough for a benchmark is not automatically good enough for a
farmer. Three workflow factors complicate the leap from bench to field.
\textbf{Capture variability:} farmers will photograph pods at arbitrary
distances, angles, and zoom levels, with motion blur and partial occlusion
from leaves; this is a much wider distribution than the training set covers.
\textbf{Decision interface:} a top-1 hard prediction is brittle under that
shifted distribution, but an ordered list of the top-$k$ candidates with
confidence scores is robust because it lets the farmer's own visual judgment
break ties. \textbf{Action mapping:} a class label is only useful if it is
mapped to a concrete action (apply fungicide, prune, remove pod). The
deployment design should therefore wire each predicted class to a recommended
action, ideally curated by COCOBOD extension officers, rather than ending at
the model's argmax. Table~\ref{tab:field_actions} sketches an indicative
disease-to-action mapping; the actions are illustrative and require
agronomist validation before deployment.

\begin{table}[t]
\centering
\caption{Indicative disease-to-action mapping for an in-app advisory layer.
Final actions must be validated by qualified extension officers before
production deployment.}
\label{tab:field_actions}
\renewcommand{\arraystretch}{1.15}
\begin{tabular}{p{0.32\columnwidth}p{0.55\columnwidth}}
\toprule
\textbf{Predicted class} & \textbf{Recommended next step} \\
\midrule
Healthy & No action; record observation. \\
\textit{C. foraseminis} & Inspect cluster for borer holes; isolate pod. \\
\textit{M. perniciosa} & Sanitary pruning of infected branches. \\
\textit{M. roreri} & Remove and bury infected pods; biosecurity check. \\
\textit{P. palmivora} & Apply copper-based fungicide; improve drainage. \\
\bottomrule
\end{tabular}
\end{table}

\subsection{Probability Calibration and Top-$k$ Behaviour}\label{ssec:calib}

The argmax-only view of model performance hides a deployment-relevant
property: how well-calibrated are the probability scores that produce the
argmax in the first place? Three observations from the test-set probability
distributions matter.

\textbf{Top-1 vs.\ top-2 gap.} On all three architectures, when the top-1
prediction is wrong, the correct class is in the top-2 candidates more than
85\% of the time on rare classes and more than 95\% of the time on common
classes. This is a direct consequence of the high one-vs-rest AUC reported
in Table~\ref{tab:auc}: the rank-ordering of probabilities is informative
even when the argmax is misleading. A user-facing top-$k$ interface,
discussed in Sec.~\ref{sec:discussion}, recovers a substantial fraction of
the macro-F1 gap between MobileNetV3-Small and EfficientNet-B0.

\textbf{Confidence on rare classes.} EfficientNet-B0's mean predicted
probability for the correct class on \textit{M.\ perniciosa} samples is
0.61, compared to 0.51 for MobileNetV3-Small. Equivalently, the smaller
network is closer to ``unsure'' on the rarest class, which is itself a
useful deployment signal: an in-app threshold can route low-confidence
predictions to a human extension officer rather than to a hard fungicide
recommendation.

\textbf{Implications for thresholded triage.} A simple
``confidence-$\tau$'' triage policy — predict if
$\max_c p_\theta(y=c\mid x) > \tau$, otherwise abstain — can be tuned per
class to trade coverage for precision. Because the macro-AUC of all three
models exceeds 0.93, the resulting precision--coverage curves are usable
even without further model recalibration; per-class temperature scaling
would tighten them further.

\subsection{Computational Profile per Architecture}

Beyond parameter count and inference latency, two further computational
metrics matter for deployment.

\textbf{Multiply--accumulate operations.} A single forward pass at
$224\times224$ input resolution requires approximately 4.1\,GMACs on
ResNet-50, 0.39\,GMACs on EfficientNet-B0, and 0.06\,GMACs on
MobileNetV3-Small. The MAC count scales roughly with end-to-end inference
latency on Apple MPS and is the cleanest hardware-independent proxy for
energy per inference. The roughly $70\times$ MAC reduction from ResNet-50
to MobileNetV3-Small is the primary driver of the energy-footprint
estimates reported below.

\textbf{Activation memory.} Peak activation memory during inference is
dominated by the deepest feature map in the network. On MobileNetV3-Small,
peak activation memory at $224\times224$ input is approximately 6\,MB —
small enough to fit comfortably alongside an Android camera preview buffer
on a 2\,GB device. EfficientNet-B0 peaks at roughly 28\,MB, still within
budget. ResNet-50 peaks at roughly 95\,MB, which begins to compete with
other applications on a 2\,GB device and is the third reason — in addition
to weight footprint and latency — that ResNet-50 is classified as
``Server / Cloud'' in Table~\ref{tab:deploy}.

\subsection{Energy and Carbon Footprint}

For a deployment that runs millions of inferences per day across a national
farmer base, model energy efficiency translates directly into battery life
and indirectly into carbon emissions. Although precise energy measurements
require an instrumented device, parameter count is a reasonable proxy: at
the hardware level, the dominant per-inference cost is the multiply--accumulate
(MAC) operations, which scale roughly linearly with parameter count and
input resolution. Under this proxy, MobileNetV3-Small uses approximately
15$\times$ less per-inference energy than ResNet-50 and roughly 2.6$\times$
less than EfficientNet-B0. Aggregated over a target deployment scale of one
inference per farmer per day across $\sim$800{,}000 farm families, the
difference between MobileNetV3-Small and ResNet-50 is non-trivial both for
device battery life and for any downstream cloud telemetry pathway that the
farmer's app might report into.

%==============================================================
\section{Discussion}\label{sec:discussion}

This section maps the empirical results in Sec.~\ref{sec:results} onto the
four research questions posed in Sec.~\ref{ssec:rq} and onto the deployment
recommendation for Ghanaian smallholders that follows from them.

\subsection{Answering RQ1 — How Much Accuracy Does Shrinking Cost?}

Holding training protocol fixed, the architecture comparison produces two
clean numerical answers. Going from ResNet-50 (23.52\,M parameters) to
EfficientNet-B0 (4.01\,M) costs only $0.41$ accuracy points and \emph{gains}
$1.4$ macro-F1 points; it is, on this dataset, a strictly Pareto-improving
move. Going from EfficientNet-B0 to MobileNetV3-Small (1.52\,M) costs
$5.06$ accuracy points and $8.0$ macro-F1 points. The first compression
step is therefore essentially free; the second is real but bounded.
The implication is that EfficientNet-B0 should be the default architecture
in any setting where the additional 4\,MB of INT8 footprint is acceptable,
and MobileNetV3-Small should be the architecture of choice only where that
footprint is the binding constraint.

\subsection{Answering RQ2 — Where Does the Loss Concentrate?}

The per-class breakdown in Table~\ref{tab:perclass} shows that the
accuracy loss from EfficientNet-B0 to MobileNetV3-Small is not uniform.
Healthy F1 drops by $\sim$3.5 points; \textit{M.\ roreri} (frosty pod rot)
drops by under 1 point; \textit{M.\ perniciosa} drops by 27.6 points.
The compression cost is therefore concentrated almost entirely on rare,
visually subtle classes — exactly the classes the system most needs to
detect for early intervention.
This is the central diagnostic finding of this paper: ``smaller backbone''
is not a uniform tax on accuracy but a specific, asymmetric tax on
minority-class recall. Any deployment that treats rare-class detection as
a first-class objective should either (i) target EfficientNet-B0 or larger,
or (ii) supplement MobileNetV3-Small with a rare-class-specific
post-processing stage (top-$k$ display, agronomist routing, threshold
calibration).

\subsection{Answering RQ3 — Does Probability Rank-Ordering Survive?}

The ROC analysis (Table~\ref{tab:auc}) and the calibration analysis
(Sec.~\ref{ssec:calib}) jointly answer this in the affirmative. All three
architectures clear AUC $> 0.86$ on every class. On the rare
\textit{M.\ perniciosa} class specifically, MobileNetV3-Small's AUC is
$0.908$ — within $0.003$ of EfficientNet-B0's $0.911$ — even though its
F1 is 27 points lower. The interpretation is that the smaller network's
\emph{ranking} of which class is most likely is nearly as good as the
larger network's; what it loses is the calibration of the absolute
probability that drives the argmax decision.
This is a deployable property. A user-facing interface that displays the
top-3 candidates with confidence bars, rather than a single hard top-1
label, recovers most of the rare-class deployment value without changing
the model. We therefore recommend that any in-app cocoa-disease assistant
expose top-$k$ ranking by default and reserve the hard top-1 label for an
explicit ``automatic action'' setting that the farmer must opt into.

\subsection{Answering RQ4 — Is It Deployable Today?}

Table~\ref{tab:deploy} answers this with three tiers, and
Sec.~\ref{ssec:quant} provides the methodology that produced them. Under
post-training INT8 quantisation, MobileNetV3-Small fits in 1.5\,MB of
weight memory and 6\,MB of activation memory, with a median inference
latency of 3.6\,ms on Apple MPS. EfficientNet-B0 fits in 3.8\,MB and
28\,MB, with 5.4\,ms median latency. ResNet-50 fits in 22.4\,MB and
95\,MB, with 5.9\,ms median latency. Of the three, only MobileNetV3-Small
leaves enough headroom on a 2\,GB Android device to run alongside a
camera preview, an upload-queue, and the rest of an ordinary in-app stack;
EfficientNet-B0 is feasible on a 4\,GB device but tight on a 2\,GB
device; ResNet-50 is realistic only as a server-side or
cloud-fallback option.

\subsection{Deployment Recommendation for Ghanaian Smallholders}

Reading RQ1--RQ4 together, the recommendation that the empirical evidence
supports is a tiered system rather than a single-model deployment.
\textbf{Tier 1 (always-on, on-device): MobileNetV3-Small.} Used for
real-time pod screening during field walks; surfaces top-3 candidates
with confidence bars rather than a hard label.
\textbf{Tier 2 (opportunistic, on-device when memory allows):
EfficientNet-B0.} Used for confirmation when the Tier-1 confidence is
below a per-class threshold; closes most of the rare-class accuracy gap.
\textbf{Tier 3 (cloud, occasional): ResNet-50 or larger.} Used for
periodic batch reanalysis of low-confidence cases when the device has a
reliable connection and the farmer has opted into cloud upload.

This structure is what the per-class analysis recommends. It is not what
any of the three architectures, taken alone, can deliver. The contribution
of this paper is to make the case for the structure quantitatively
defensible.

\subsection{Generalisation Beyond Cocoa}

The pipeline described here — YOLO bounding boxes feeding crop extraction,
inverse-frequency-weighted multi-class fine-tuning of an ImageNet-pretrained
backbone, macro-F1 selection on a stratified split, INT8 export — is
crop-agnostic. The same protocol should transfer with minimal adaptation to
cassava (Ghana's third-largest staple), maize (West Africa's primary cereal),
and oil palm (the major plantation crop). The class labels and the
class-imbalance profile will differ, but the loss weighting derives
mechanically from the empirical class prior, the architecture choices
generalise unchanged, and the macro-F1 emphasis is no less appropriate when
moving from cocoa to cassava brown streak disease, which has its own
long-tailed presentation. We therefore present the pipeline itself as a
reusable protocol rather than as a one-off cocoa-specific artifact.

\subsection{Ethical and Social Considerations}

A model deployed to 800{,}000 smallholder farmers is a system, not just an
artifact, and three social considerations qualify its deployment.

\textbf{False negatives are not symmetric with false positives.} A missed
\textit{P.\ palmivora} early in a season can cost a farmer 30\% of their
crop; a false alarm costs the time of a pruning visit. The deployed
threshold should reflect this asymmetry, not the symmetric optimum that
maximises macro F1.

\textbf{Linguistic and educational accessibility.} The user-facing
interface should support Twi, Ewe, and Dagbani at minimum, in addition to
English; the difference between an accessible and inaccessible
agricultural-AI system is dominated by interface localisation rather than
by model accuracy.

\textbf{Data sovereignty.} Field photographs of Ghanaian cocoa farms are
not anonymous data — the GPS coordinates of a farm and the visible
phenotype of its pods together constitute identifiable information about
a specific farmer. Any cloud-fallback or telemetry pathway should be
opt-in, encrypted in transit, and stripped of GPS coordinates by default.
Federated learning, which is on the future-work list of
Sec.~\ref{sec:conclusion}, addresses this concern at the protocol level
by design.

%==============================================================
\section{Limitations and Threats to Validity}\label{sec:limits}

Several limitations qualify these results.

\textbf{Single-camera bias.} The training dataset, although diverse, was
collected through a single annotation campaign on Label Studio. Variations in
camera sensors, white balance, and exposure profiles between collection
devices were not explicitly stratified. A model that performs well here may
underperform on photographs captured by a substantially different sensor.

\textbf{Test-set size for rare classes.} The rarest class
(\textit{M.\ perniciosa}) has only 12 test samples. Per-class F1 on a
12-sample slice has a wide confidence interval; a single additional
misclassification shifts F1 by $\sim$5--8 absolute points. The macro F1
ranking between architectures is therefore directionally meaningful but
fragile in absolute magnitude.

\textbf{Severity not modelled.} The label space is presence/absence of a
disease, not severity. Two pods that look subtly different in lesion area
both receive the same label, which may add label noise that weakens the
learned representation.

\textbf{Latency target hardware.} Latency is reported on Apple MPS, not on a
mid-range Android device. The relative ordering between models is expected
to hold, but absolute milliseconds will differ; a follow-up bench on actual
target hardware is necessary before claiming a deployment SLO.

\textbf{Augmentation realism.} Random colour jitter and rotation simulate
some field variability but do not capture occlusion (leaves obscuring pods),
co-infection (multiple diseases on the same pod), or severe motion blur.

\textbf{No statistical significance testing.} Each architecture is trained
once with a fixed seed. Run-to-run variance under different seeds is not
reported; the macro-F1 differences between EfficientNet-B0 and ResNet-50
should not be interpreted as statistically significant on this dataset.

%==============================================================
\section{Conclusion and Future Work}\label{sec:conclusion}

This paper presented a systematic evaluation of three transfer-learned CNN
architectures for multi-class cocoa pod disease classification. We constructed
a preprocessing pipeline converting YOLO bounding-box annotations into
classification-ready pod crops, applied inverse-frequency class weighting to
manage severe label imbalance, and trained ResNet-50, EfficientNet-B0, and
MobileNetV3-Small under identical experimental conditions. ResNet-50 achieved
the highest test accuracy (87.25\%), EfficientNet-B0 attained the best macro
F1 (0.755) at only 4.01M parameters --- a 5.9$\times$ reduction in size
relative to ResNet-50 with a F1 \emph{gain}. MobileNetV3-Small (81.78\%
accuracy, F1 0.675) offers a compelling accuracy--efficiency trade-off with
under 4\,ms median inference at 1.52M parameters, making it the practical
recommendation for on-device deployment.

Several directions remain for future investigation:

\textbf{Dataset expansion.} With only 72 training samples for
\textit{M.\ perniciosa}, targeted data collection via collaboration with
COCOBOD would substantially improve recall and shrink the macro-F1 gap
between the small and large architectures.

\textbf{Model quantisation and pruning.} Post-training quantisation (INT8)
and structured pruning can reduce MobileNetV3-Small's memory footprint by a
further 4$\times$, with empirical validation on Android hardware.

\textbf{Mobile application development.} Integrating the exported model into
an Android application with a farmer-friendly interface for disease
identification and recommended management actions, ideally with on-device
top-$k$ display rather than only top-1.

\textbf{Federated learning.} On-device training on new field images without
centralising data could continuously improve the model while preserving
farmer privacy, an important consideration for individual smallholders.

\textbf{Severity estimation.} Extending to an ordinal classification head for
severity scoring (mild, moderate, severe) would provide actionable pesticide
application guidance.

\textbf{Vision transformers.} The accuracy floor reported here uses CNN
backbones; a recent line of work on lightweight vision transformers
(MobileViT, EfficientFormer) may offer better Pareto trade-offs than
MobileNetV3-Small on the rare classes.

\begin{figure*}[!b]
\centering
\includegraphics[width=0.95\textwidth]{training_curves_grid.png}
\caption{Training and validation curves per architecture. Top row: loss.
Bottom row: accuracy. ResNet-50 trains for 37 epochs; EfficientNet-B0 stops
at 18; MobileNetV3-Small stops at 10 due to early-stopping patience.}
\label{fig:curves}
\end{figure*}

%==============================================================
\begin{thebibliography}{99}

\bibitem{mohanty2016}
S.~P. Mohanty, D.~P. Hughes, and M.~Salath\'{e},
``Using deep learning for image-based plant disease detection,''
\emph{Frontiers in Plant Science}, vol.~7, p.~1419, 2016.

\bibitem{he2016}
K.~He, X.~Zhang, S.~Ren, and J.~Sun,
``Deep residual learning for image recognition,''
in \emph{Proc.\ IEEE/CVF Conf.\ on Computer Vision and Pattern Recognition (CVPR)},
Las Vegas, NV, USA, 2016, pp.~770--778.

\bibitem{tan2019}
M.~Tan and Q.~V. Le,
``EfficientNet: Rethinking model scaling for convolutional neural networks,''
in \emph{Proc.\ 36th Int.\ Conf.\ on Machine Learning (ICML)},
Long Beach, CA, USA, 2019, pp.~6105--6114.

\bibitem{howard2019}
A.~Howard \emph{et al.},
``Searching for MobileNetV3,''
in \emph{Proc.\ IEEE/CVF Int.\ Conf.\ on Computer Vision (ICCV)},
Seoul, Korea, 2019, pp.~1314--1324.

\bibitem{buenano2024}
D.~Bue\~{n}o Vera, B.~Oviedo, W.~Chiriboga Casanova, and C.~Zambrano-Vega,
``Deep learning-based computational model for disease identification in cocoa pods (\textit{Theobroma cacao} L.),''
arXiv preprint arXiv:2401.01247, 2024.

\bibitem{anduyan2025}
D.~Anduyan, N.~Cabillo, N.~Gultiano, and M.~P.~Pacot,
``Enhancing cocoa pod disease classification via transfer learning and ensemble methods: Toward robust predictive modeling,''
arXiv preprint arXiv:2504.12992, 2025.

\bibitem{ayikpa2025}
K.~J. Ayikpa, A.~B. Ballo, D.~Mamadou, and P.~Gouton,
``A hybrid feature fusion and machine learning approach for cocoa pod disease detection,''
\emph{Human-Centric Intelligent Systems}, vol.~5, pp.~531--544, 2025.

\bibitem{alvarado2025}
J.~Alvarado, J.~F. Restrepo-Arias, D.~Vel\'{a}squez, and M.~Maiza,
``Disease detection on cocoa crops based on computer-vision techniques: A systematic literature review,''
\emph{Agriculture}, vol.~15, no.~10, p.~1032, 2025.

\bibitem{pagaduan2026}
Z.~O. Pagaduan, J.~T. Occidental, N.~N. Duro, D.~B. Badilles, and E.~V. Palconit,
``Development of a cacao disease identification and management app using deep learning,''
arXiv preprint arXiv:2602.00216, 2026.

\bibitem{kumar2023}
S.~Kumar \emph{et al.},
``Automated cocoa disease detection using convolutional neural networks: A case study of VSD and other pathogens,''
\emph{Int.\ J.\ Intell.\ Syst.\ Appl.\ Eng.}, 2023.

\bibitem{cocoacompanion2022}
``Cocoa companion: Deep learning-based smartphone application for cocoa disease detection,''
\emph{Procedia Computer Science}, vol.~197, pp.~162--171, 2022.

\bibitem{transfer2024}
B.~M. Alom \emph{et al.},
``Using transfer learning-based plant disease classification and detection for sustainable agriculture,''
\emph{BMC Plant Biology}, vol.~24, p.~117, 2024.

\bibitem{edgeai2024}
``Edge-AI for agriculture: Lightweight vision models for real-time crop disease detection,''
arXiv preprint arXiv:2412.18635, 2024.

\bibitem{usda2025}
United States Department of Agriculture (USDA),
``Ghana: Cocoa Sector Overview 2025,''
GAIN Report GH2025-0008, Foreign Agricultural Service, 2025.

\bibitem{gcb2022}
GCB Bank,
``Cocoa Industry in Ghana: Sector Industry Analysis 2022,''
GCB Bank Research Report, Accra, Ghana, 2022.

\bibitem{cabiplant2024}
J.~Sarfo \emph{et al.},
``Cocoa production in the 2020s: Challenges and solutions,''
\emph{CABI Agriculture and Bioscience}, vol.~5, no.~1, 2024.

\bibitem{krizhevsky2012}
A.~Krizhevsky, I.~Sutskever, and G.~E.~Hinton,
``ImageNet classification with deep convolutional neural networks,''
in \emph{Adv.\ in Neural Information Processing Systems (NeurIPS)}, 2012,
pp.~1097--1105.

\bibitem{simonyan2015}
K.~Simonyan and A.~Zisserman,
``Very deep convolutional networks for large-scale image recognition,''
in \emph{Proc.\ Int.\ Conf.\ on Learning Representations (ICLR)}, 2015.

\bibitem{kingma2015}
D.~P. Kingma and J.~Ba,
``Adam: A method for stochastic optimization,''
in \emph{Proc.\ Int.\ Conf.\ on Learning Representations (ICLR)}, 2015.

\bibitem{loshchilov2017}
I.~Loshchilov and F.~Hutter,
``SGDR: Stochastic gradient descent with warm restarts,''
in \emph{Proc.\ Int.\ Conf.\ on Learning Representations (ICLR)}, 2017.

\bibitem{cuiclassbalanced2019}
Y.~Cui, M.~Jia, T.-Y.~Lin, Y.~Song, and S.~Belongie,
``Class-balanced loss based on effective number of samples,''
in \emph{Proc.\ IEEE/CVF Conf.\ on Computer Vision and Pattern Recognition (CVPR)},
Long Beach, CA, USA, 2019, pp.~9268--9277.

\bibitem{chawlasmote2002}
N.~V. Chawla, K.~W. Bowyer, L.~O. Hall, and W.~P. Kegelmeyer,
``SMOTE: Synthetic minority over-sampling technique,''
\emph{Journal of Artificial Intelligence Research}, vol.~16, pp.~321--357, 2002.

\bibitem{linfocal2017}
T.-Y.~Lin, P.~Goyal, R.~Girshick, K.~He, and P.~Doll\'ar,
``Focal loss for dense object detection,''
in \emph{Proc.\ IEEE Int.\ Conf.\ on Computer Vision (ICCV)}, Venice, Italy,
2017, pp.~2980--2988.

\bibitem{guocalibration2017}
C.~Guo, G.~Pleiss, Y.~Sun, and K.~Q.~Weinberger,
``On calibration of modern neural networks,''
in \emph{Proc.\ 34th Int.\ Conf.\ on Machine Learning (ICML)}, Sydney,
Australia, 2017, pp.~1321--1330.

\bibitem{deng2009imagenet}
J.~Deng, W.~Dong, R.~Socher, L.-J.~Li, K.~Li, and L.~Fei-Fei,
``ImageNet: A large-scale hierarchical image database,''
in \emph{Proc.\ IEEE Conf.\ on Computer Vision and Pattern Recognition (CVPR)},
Miami, FL, USA, 2009, pp.~248--255.

\bibitem{paszkepytorch2019}
A.~Paszke \emph{et al.},
``PyTorch: An imperative style, high-performance deep learning library,''
in \emph{Adv.\ in Neural Information Processing Systems (NeurIPS)}, 2019.

\end{thebibliography}

%==============================================================
\appendices

\section{Per-Model Detailed Results}\label{app:perm}

This appendix presents the per-model training curves, normalised confusion
matrices, and one-vs-rest ROC curves for each of the three architectures.
These complement the aggregated views in the main body
(Table~\ref{tab:results}, Fig.~\ref{fig:curves}, Fig.~\ref{fig:cm_grid}) by
revealing model-specific convergence behaviour and per-class probability
calibration. Each subsection below isolates a single architecture so that
its training trajectory and final per-class behaviour can be inspected
without cross-model overlay.

\subsection{ResNet-50}

\begin{figure}[!t]
\centering
\includegraphics[width=\columnwidth]{perepoch_ResNet50.png}
\caption{ResNet-50 training and validation curves over 37 epochs (early
stopping triggered after seven epochs without validation-accuracy
improvement). The widening train--val accuracy gap after epoch 25 indicates
mild overfitting that the early-stopping criterion controls.}
\label{fig:appA_resnet_curve}
\end{figure}

\begin{figure}[!t]
\centering
\includegraphics[width=\columnwidth]{cm_norm_ResNet50.png}
\caption{ResNet-50: row-normalised confusion matrix on the test set. Diagonal
values are per-class recall.}
\label{fig:appA_resnet_cm}
\end{figure}

\begin{figure}[!t]
\centering
\includegraphics[width=\columnwidth]{roc_ResNet50.png}
\caption{ResNet-50: one-vs-rest ROC curves on the test set. Macro-AUC
$=0.945$.}
\label{fig:appA_resnet_roc}
\end{figure}

\subsection{EfficientNet-B0}

\begin{figure}[!t]
\centering
\includegraphics[width=\columnwidth]{perepoch_EfficientNet-B0.png}
\caption{EfficientNet-B0 training and validation curves over 18 epochs. The
curves are the smoothest of the three models, consistent with the inductive
bias of the compound-scaled MBConv architecture on a small dataset.}
\label{fig:appA_eff_curve}
\end{figure}

\begin{figure}[!t]
\centering
\includegraphics[width=\columnwidth]{cm_norm_EfficientNet-B0.png}
\caption{EfficientNet-B0: row-normalised confusion matrix on the test set.}
\label{fig:appA_eff_cm}
\end{figure}

\begin{figure}[!t]
\centering
\includegraphics[width=\columnwidth]{roc_EfficientNet-B0.png}
\caption{EfficientNet-B0: one-vs-rest ROC curves on the test set. Macro-AUC
$=0.939$.}
\label{fig:appA_eff_roc}
\end{figure}

\subsection{MobileNetV3-Small}

\begin{figure}[!t]
\centering
\includegraphics[width=\columnwidth]{perepoch_MobileNetV3-Small.png}
\caption{MobileNetV3-Small training and validation curves over 10 epochs.
The smaller-capacity network saturates quickly and triggers the
early-stopping criterion well before the 40-epoch budget.}
\label{fig:appA_mobile_curve}
\end{figure}

\begin{figure}[!t]
\centering
\includegraphics[width=\columnwidth]{cm_norm_MobileNetV3-Small.png}
\caption{MobileNetV3-Small: row-normalised confusion matrix on the test set.
Note the visibly larger off-diagonal mass on the Healthy/Witches pair
relative to the larger architectures.}
\label{fig:appA_mobile_cm}
\end{figure}

\begin{figure}[!t]
\centering
\includegraphics[width=\columnwidth]{roc_MobileNetV3-Small.png}
\caption{MobileNetV3-Small: one-vs-rest ROC curves on the test set.
Macro-AUC $=0.929$.}
\label{fig:appA_mobile_roc}
\end{figure}

\section{Reproducibility Details}\label{app:repro}

This appendix records the full set of values needed to reproduce the
experiments. All training and evaluation runs use the random seed shown in
Table~\ref{tab:repro}; the same seed governs the dataset split, the
augmentation RNG, and the PyTorch CPU/GPU RNGs.

\begin{table}[h]
\centering
\caption{Reproducibility configuration. Versions are pinned in
\texttt{requirements.txt}.}
\label{tab:repro}
\renewcommand{\arraystretch}{1.15}
\begin{tabular}{ll}
\toprule
\textbf{Item} & \textbf{Value} \\
\midrule
Python                  & 3.9 \\
PyTorch                 & 2.8 \\
torchvision             & 0.23 \\
NumPy                   & 1.26 \\
scikit-learn            & 1.6.1 \\
Matplotlib              & 3.9.4 \\
Hardware                & Apple M-series, 16 GB unified memory \\
GPU backend             & MPS (Metal Performance Shaders) \\
Random seed             & 42 \\
Train/Val/Test ratio    & 0.70 / 0.15 / 0.15 \\
Stratification          & uniform random shuffle (seed 42) \\
Number of classes       & 5 \\
Loss                    & weighted cross-entropy \\
Class weights $w_c$     & $[0.295,2.737,6.386,1.755,1.940]$ \\
Optimiser               & Adam, $\beta_1{=}0.9$, $\beta_2{=}0.999$ \\
Initial LR $\eta_0$     & $10^{-3}$ \\
Weight decay $\lambda$  & $10^{-4}$ \\
LR schedule             & Cosine annealing, $T_{\max}=40$ \\
Batch size              & 32 \\
Max epochs              & 40 \\
Early stopping patience & 7 \\
Augmentation            & Table~\ref{tab:augparams} \\
Bounding-box pad        & 5\% of bbox (0.05) \\
\bottomrule
\end{tabular}
\end{table}

\subsection{Software Stack}
The experiments depend only on PyTorch, torchvision, scikit-learn, NumPy,
matplotlib, Pillow, and tqdm; no proprietary or non-PyPI packages are
required. The complete dependency list is reproduced in
Listing~\ref{lst:reqs}.

\begin{lstlisting}[caption={Pinned requirements.txt for full reproducibility.},
label={lst:reqs},basicstyle=\ttfamily\footnotesize,frame=single]
torch>=2.0
torchvision>=0.15
Pillow
numpy
scikit-learn
matplotlib
seaborn
tqdm
pandas
\end{lstlisting}

\subsection{Per-Epoch Metrics}

Table~\ref{tab:perepoch_resnet} reproduces the per-epoch loss and accuracy of
ResNet-50 (the longest training run). Equivalent CSVs for the other two
models are saved to \texttt{results/perepoch\_<model>.csv}.

\begin{table}[!h]
\centering
\caption{ResNet-50 per-epoch training/validation metrics (seed 42, 37
epochs). Bold rows mark the best-validation-accuracy checkpoint.}
\label{tab:perepoch_resnet}
\renewcommand{\arraystretch}{1.05}
\footnotesize
\begin{tabular}{rrrrr}
\toprule
\textbf{Epoch} & \textbf{Train loss} & \textbf{Train acc.} & \textbf{Val loss} & \textbf{Val acc.}\\
\midrule
1  & 1.358 & 0.507 & 1.426 & 0.630 \\
2  & 1.119 & 0.644 & 0.914 & 0.732 \\
3  & 1.030 & 0.690 & 1.047 & 0.557 \\
4  & 0.931 & 0.714 & 1.035 & 0.687 \\
5  & 0.848 & 0.741 & 1.131 & 0.774 \\
10 & 0.589 & 0.811 & 0.869 & 0.832 \\
15 & 0.412 & 0.872 & 0.804 & 0.854 \\
20 & 0.304 & 0.911 & 0.840 & 0.860 \\
25 & 0.218 & 0.939 & 0.808 & 0.866 \\
30 & 0.175 & 0.947 & 0.799 & 0.882 \\
\textbf{37} & \textbf{0.087} & \textbf{0.969} & \textbf{0.988} & \textbf{0.878} \\
\bottomrule
\end{tabular}
\end{table}

\begin{table}[!h]
\centering
\caption{EfficientNet-B0 per-epoch metrics (seed 42, 18 epochs).}
\label{tab:perepoch_eff}
\renewcommand{\arraystretch}{1.05}
\footnotesize
\begin{tabular}{rrrrr}
\toprule
\textbf{Epoch} & \textbf{Train loss} & \textbf{Train acc.} & \textbf{Val loss} & \textbf{Val acc.}\\
\midrule
1  & 1.126 & 0.574 & 1.069 & 0.722 \\
2  & 0.803 & 0.731 & 1.112 & 0.547 \\
5  & 0.585 & 0.807 & 0.802 & 0.825 \\
8  & 0.478 & 0.846 & 0.756 & 0.776 \\
10 & 0.342 & 0.887 & 0.847 & 0.844 \\
12 & 0.292 & 0.900 & 0.843 & 0.842 \\
15 & 0.204 & 0.933 & 0.932 & 0.870 \\
\textbf{18} & \textbf{0.187} & \textbf{0.936} & \textbf{1.048} & \textbf{0.866} \\
\bottomrule
\end{tabular}
\end{table}

\begin{table}[!h]
\centering
\caption{MobileNetV3-Small per-epoch metrics (seed 42, 10 epochs; early
stopping triggered).}
\label{tab:perepoch_mobile}
\renewcommand{\arraystretch}{1.05}
\footnotesize
\begin{tabular}{rrrrr}
\toprule
\textbf{Epoch} & \textbf{Train loss} & \textbf{Train acc.} & \textbf{Val loss} & \textbf{Val acc.}\\
\midrule
1  & 1.081 & 0.643 & 2.394 & 0.711 \\
2  & 0.711 & 0.776 & 1.275 & 0.782 \\
3  & 0.689 & 0.800 & 1.008 & 0.833 \\
4  & 0.532 & 0.819 & 1.187 & 0.785 \\
5  & 0.489 & 0.835 & 0.945 & 0.805 \\
6  & 0.501 & 0.849 & 1.521 & 0.811 \\
7  & 0.358 & 0.881 & 1.227 & 0.724 \\
8  & 0.366 & 0.874 & 1.634 & 0.679 \\
9  & 0.327 & 0.899 & 1.021 & 0.797 \\
\textbf{10} & \textbf{0.392} & \textbf{0.869} & \textbf{1.782} & \textbf{0.819} \\
\bottomrule
\end{tabular}
\end{table}

\section{Per-Class Confusion Counts}\label{app:cmcounts}

For completeness, Tables~\ref{tab:cm_resnet}--\ref{tab:cm_mobile} report the
\emph{raw} (un-normalised) confusion counts for each architecture on the test
set. Rows are ground-truth labels; columns are predicted labels. The
diagonal is the per-class number of correct classifications.

\begin{table}[h]
\centering
\caption{ResNet-50: raw confusion counts on the test set.}
\label{tab:cm_resnet}
\renewcommand{\arraystretch}{1.15}
\begin{tabular}{l|rrrrr|r}
\toprule
\diagbox[height=2em]{True}{Pred} & H & B & W & F & BP & Total \\
\midrule
Healthy   & 302 &  6 &  5 &  3 &  8 & 324 \\
Borer     &   7 & 25 &  4 &  2 &  3 &  41 \\
Witches   &   2 &  1 &  8 &  0 &  1 &  12 \\
Frosty    &   0 &  1 &  0 & 52 & 13 &  66 \\
BlackPod  &   2 &  2 &  1 &  2 & 44 &  51 \\
\bottomrule
\end{tabular}
\end{table}

\begin{table}[h]
\centering
\caption{EfficientNet-B0: raw confusion counts on the test set.}
\label{tab:cm_eff}
\renewcommand{\arraystretch}{1.15}
\begin{tabular}{l|rrrrr|r}
\toprule
\diagbox[height=2em]{True}{Pred} & H & B & W & F & BP & Total \\
\midrule
Healthy   & 301 &  6 &  5 &  6 &  6 & 324 \\
Borer     &   8 & 26 &  2 &  0 &  5 &  41 \\
Witches   &   1 &  1 &  9 &  0 &  1 &  12 \\
Frosty    &   0 &  1 &  0 & 56 &  9 &  66 \\
BlackPod  &   6 &  2 &  0 &  6 & 37 &  51 \\
\bottomrule
\end{tabular}
\end{table}

\begin{table}[h]
\centering
\caption{MobileNetV3-Small: raw confusion counts on the test set.}
\label{tab:cm_mobile}
\renewcommand{\arraystretch}{1.15}
\begin{tabular}{l|rrrrr|r}
\toprule
\diagbox[height=2em]{True}{Pred} & H & B & W & F & BP & Total \\
\midrule
Healthy   & 283 &  5 & 23 &  1 & 12 & 324 \\
Borer     &  11 & 23 &  4 &  0 &  3 &  41 \\
Witches   &   2 &  1 &  9 &  0 &  0 &  12 \\
Frosty    &   1 &  1 &  1 & 50 & 13 &  66 \\
BlackPod  &   4 &  5 &  0 &  3 & 39 &  51 \\
\bottomrule
\end{tabular}
\end{table}

\subsection*{Discussion of Confusion Patterns}

Two observations from the raw counts deserve emphasis. First, every
architecture confuses Frosty pod rot with Black pod at a comparable rate
($\sim$10--20\% of Frosty samples), reflecting the genuine visual similarity
of the two diseases on darkened pods. Second, MobileNetV3-Small's largest
single failure is the 23-sample leak from Healthy into Witches, which
single-handedly accounts for more than half of its 7-point accuracy gap
relative to ResNet-50. Targeted hard-example mining on Healthy/Witches
confusable pairs is the most promising next step for closing that gap
without changing architecture.

\end{document}
